% !TEX program = xelatex
% 由【完整版】马原期末复习.md 忠实转换而来（结构 1:1，含 6 层嵌套；正文未改动、未扩写）。
\documentclass[12pt,openany]{ctexbook}

\usepackage{amsmath,amssymb}
\usepackage{xcolor}
\usepackage{graphicx}
\usepackage{enumitem}
\usepackage{array,booktabs,tabularx}
\usepackage[most]{tcolorbox}
\usepackage[a4paper,margin=2.4cm]{geometry}
\usepackage{hyperref}
\hypersetup{hidelinks,bookmarksnumbered=false,
  pdftitle={马克思主义基本原理（知识点梳理）},pdfauthor={周睿}}

% 标题不自动编号：原文标题已含「第一章」等
\setcounter{secnumdepth}{-1}
\setcounter{tocdepth}{1}

% markdown *强调* 用楷体（非斜体，文本层可复制可检索）
\newcommand{\zhemph}[1]{{\kaishu #1}}

% —— 深层列表支持：itemize/enumerate 均放开到 9 层（本文用到 6 层）——
\setlistdepth{9}
\renewlist{itemize}{itemize}{9}
\renewlist{enumerate}{enumerate}{9}
\setlist[itemize]{leftmargin=1.5em,itemsep=1.5pt,topsep=2pt,parsep=0pt}
\setlist[itemize,1]{label=$\bullet$}
\setlist[itemize,2]{label=$\circ$}
\setlist[itemize,3]{label=\textbf{--}}
\setlist[itemize,4]{label=$\ast$}
\setlist[itemize,5]{label=$\cdot$}
\setlist[itemize,6]{label=$\triangleright$}
\setlist[itemize,7]{label=$\bullet$}
\setlist[itemize,8]{label=$\circ$}
\setlist[itemize,9]{label=\textbf{--}}
\setlist[enumerate]{leftmargin=2em,itemsep=1.5pt,topsep=2pt}
\setlist[enumerate,1]{label=\arabic*.}
\setlist[enumerate,2]{label=(\alph*)}
\setlist[enumerate,3]{label=\roman*.}
\setlist[enumerate,4]{label=\Alph*.}
\setlist[enumerate,5]{label=\arabic*.}
\setlist[enumerate,6]{label=(\alph*)}
\setlist[enumerate,7]{label=\roman*.}
\setlist[enumerate,8]{label=\Alph*.}
\setlist[enumerate,9]{label=\arabic*.}

% 引用块（markdown blockquote）
\newtcolorbox{mdquote}{
  enhanced,breakable,
  colback=black!4,colframe=black!45,
  boxrule=0pt,leftrule=2.6pt,arc=0pt,outer arc=0pt,
  left=10pt,right=8pt,top=6pt,bottom=6pt,boxsep=0pt,
  before skip=8pt,after skip=8pt,
}

\setlength{\parindent}{0pt}
\setlength{\parskip}{0.45em plus 2pt minus 1pt}
\linespread{1.16}

\title{\bfseries 马克思主义基本原理\\[2pt]\large 知识点梳理}
\author{周睿}
\date{2024 年 1 月}

\begin{document}
\frontmatter
\maketitle
\tableofcontents
\mainmatter
\chapter{导论}
\section{什么是马克思主义}
\begin{itemize}
  \item 定义：
\begin{itemize}
  \item 马克思主义是由马克思恩格斯创立并为后继者所不断发展的科学理论体系；
  \item 是关于自然、社会和人类思维发展的一般规律学说；
  \item 是关于社会主义必然代替资本主义，最终实现共产主义的学说；
  \item 是关于无产阶级解放、全人类解放和每个人自由而全面发展的学说；
  \item 是无产阶级政党和社会主义国家的指导思想，是指引人民创造美好生活的行动指南。
\end{itemize}
  \item 内容：
\begin{itemize}
  \item 马克思主义的三个基本组成部分是马克思主义哲学、马克思主义政治经济学和科学社会主义。
  \item 马克思主义基本原理是对马克思主义立场、观点、方法的集中概括，是马克思主义在其形成、发展和运用过程中经过实践反复检验而确立起来的具有普遍真理性的理论。
\end{itemize}
  \item 基本立场、基本观点和基本方法：
\begin{itemize}
  \item 基本立场：马克思主义观察、分析和解决问题的根本立足点和出发点。
  \item 基本观点：关于自然、社会和人类思维发展一般规律的科学认识，是对人类思想成果和社会实践经验的科学总结。
  \item 基本方法：马克思主义的基本方法，是建立在辩证唯物主义和历史唯物主义世界观和方法论基础上，指导我们正确认识世界和改造世界的思想方法和工作方法。
\end{itemize}
\end{itemize}
\section{马克思主义的产生}
\begin{itemize}
  \item 马克思主义的创立：
\begin{itemize}
  \item 马克思主义产生于19世纪40年代，创始人是马克思和恩格斯。
  \item 马克思主义的产生有深刻的社会根源、阶级基础和思想渊源：
\begin{itemize}
  \item 社会根源：资本主义生产方式在西欧有了相当大的发展，一方面带来了社会化大生产的迅猛发展，另一方面又造成了深重的社会灾难。
  \item 阶级基础：无产阶级在反抗资产阶级剥削和压迫的斗争中，逐步走向自觉，并迫切渴望科学的理论指导。
  \item 思想渊源：
\begin{itemize}
  \item 19世纪西欧三大先进思潮，德国古典哲学、英国古典政治经济学、英法两国的空想社会主义，为马克思主义的建立提供了直接的理论来源； 
  \item 19世纪的三大科学发现，即细胞学说、能量守恒与转化定律、生物进化论，为马克思主义的产生提供了自然科学前提。
\end{itemize}
\end{itemize}
\end{itemize}
  \item 马克思主义的发展：
\begin{itemize}
  \item 列宁在领导俄国革命和建设的过程中，把马克思主义基本原理与俄国实际相结合，创立了列宁主义，把马克思主义发展到一个新的历史阶段；
  \item 中国共产党从成立起，就把马克思列宁主义确立为指导思想，并在不断探索中把马克思主义基本原理同中国具体实际相结合，领导全国各族人民取得了革命、建设、改革的伟大胜利，并不断推进马克思主义中国化，产生了毛泽东思想、邓小平理论、“三个代表”重要思想、科学发展观、习近平新时代中国特色社会主义思想，丰富和发展了马克思主义。
\end{itemize}
\end{itemize}
\section{马克思主义的鲜明特征}
\begin{itemize}
  \item 科学性：马克思主义是对自然、社会和人类思维发展本质和规律的正确反映；
  \item 人民性：人民至上是马克思主义的政治立场；
  \item 实践性：马克思主义是从实践中来、到实践中去，在实践中接受检验、并随实践而不断发展的学说；
  \item 发展性：马克思主义是不断发展的学说，具有与时俱进的理论品质。
\end{itemize}
\chapter{第一章\quad{}世界的物质性及其发展规律}
\section{哲学基本问题}
\begin{itemize}
  \item 哲学的基本问题：\textbf{思维与存在的关系问题}
\begin{itemize}
  \item \textbf{存在和思维、物质和意识谁为本原的问题，即何者为第一性的问题}
  \par →形成唯物主义和唯心主义两种根本对立的哲学派别
  \item \textbf{存在和思维、物质和意识是否具有同一性的问题，即思维能否正确地反映存在、人能否认识或彻底认识世界的问题}
  \par →产生可知论和不可知论的理论分野
  \par 对哲学基本问题的回答是解决其他一切哲学问题的前提和基础，只有科学解决存在和思维的关系问题，才能正确认识世界的本质和把握世界发展的规律。
\end{itemize}
\end{itemize}
\section{马克思主义物质观}
\subsection{物质及其存在方式}
\begin{itemize}
  \item 物质的定义：
\begin{itemize}
  \item \textbf{所谓物质，就是不依赖人的意识而存在，并能为人的意识所反映的客观存在。}
  \item 马克思主义的物质范畴从现实存在着的自然与社会存在中抽象出了其共同特性——客观实在性。
\end{itemize}
  \item 根本属性：运动
\begin{itemize}
  \item \textbf{物质的根本属性是运动。运动是标志着一切事物和现象的变化及其过程的哲学范畴。}
  \item 物质和运动是不可分割的，运动是物质的运动，物质是运动着的物质，离开物质的运动和离开运动的物质都是不可想象的。
  \item 运动的绝对性和相对性：
\begin{itemize}
  \item 物质世界的运动是绝对的，而物质在运动过程中又有某种相对的静止。
  \item 相对静止是物质运动在一定条件下的稳定状态，具体包括两种状态：
  \par 空间的相对位置的暂时不变 \& 事物的根本性质暂时不变。
  \item 运动的绝对性体现了物质运动的变动性、无条件性，静止的相对性体现了物质运动的稳定性、有条件性。
  \item 运动和静止相互依赖、相互渗透、相互包含，“动中有静、静中有动”。无条件的绝对运动和有条件的相对静止构成了对立统一的关系。
\end{itemize}
\end{itemize}
  \item 基本存在形式：时空
\begin{itemize}
  \item 时间和空间是运动着的物质的基本存在形式。
\begin{itemize}
  \item 时间是指物质运动的持续性、顺序性，特点是一维性，即时间一去不复返；
  \item 空间是指物质运动的广延性、伸张性，特点是三维性，即空间具有长宽高三方面的规定性。
\end{itemize}
  \item 物质运动总是在一定的时间和空间中发生的，没有离开物质运动的“纯粹”时间和空间，也没有离开时间和空间的物质运动。
  \item 物质运动与时间和空间的不可分割性，证明了时间和空间的客观性。
  \item 具体物质形态的时空是有限的，而整个物质世界的时空是无限的。
  \par 物质、运动、时间、空间具有内在统一性，要求我们想问题、办事情都要以具体的时间、地点和条件为转移
\end{itemize}
  \item 理论意义：
\begin{itemize}
  \item \textbf{坚持了唯物主义一元论，同唯心主义一元论和二元论划清了界限}；
  \item \textbf{坚持了能动的反映论和可知论，批判了不可知论}；
  \item \textbf{体现了唯物论和辩证法的统一，克服了形而上学唯物主义的缺陷}；
  \item 体现了\textbf{唯物主义自然观和唯物主义历史观的统一}，为彻底的唯物主义奠定了理论基础。
\end{itemize}
\end{itemize}
\subsection{物质与意识的辩证关系}
\begin{itemize}
  \item 物质决定意识：物质对意识的决定作用表现在意识的起源和本质上。
\begin{itemize}
  \item 从意识的起源来看：
\begin{itemize}
  \item \textbf{意识既是自然界长期发展的产物，也是社会历史发展的产物}。
  \item 社会实践，特别是劳动，在意识的产生和发展中起着决定性的作用。劳动为意识的产生和发展提供了客观需要和可能，在人们的劳动和交往中形成的语言促进了意识的发展。
\end{itemize}
  \item 从意识的本质来看：
\begin{itemize}
  \item \textbf{意识是人脑这样一种特殊物质的机能和属性，是客观世界的主观映象}。
  \item 因此，意识是客观内容和主观形式的统一。意识是物质的产物，但又不是物质本身。
\end{itemize}
\end{itemize}
  \item 意识对物质具有能动的反作用：
\begin{itemize}
  \item 意识具有目的性和计划性；
  \item 意识活动具有创造性：不仅会反映事物的外部现象，还会进行加工制作与选择建构；
  \item 意识具有指导实践改造客观世界的作用：以观念为指导，通过实践使之一步步成为客观现实；
  \item 意识具有调控人的行为和生理活动的作用。
\end{itemize}
\end{itemize}
\subsection{世界的物质统一性}
马克思主义认为，世界的统一性在于它的物质性，世界统一于物质。世界的物质统一性是多样性的统一。
\begin{itemize}
  \item 自然界是物质的；
  \item 人类社会本质上也是物质的；
  \item 人的意识统一于物质。
\end{itemize}
\section{马克思主义实践观}
\begin{itemize}
  \item 实践的定义：\textbf{实践是人类能动地改造世界的社会性的物质活动}。
  \item 实践的特征：
\begin{itemize}
  \item \textbf{客观实在性}：实践是人类改造世界的客观物质活动，它虽然是人类有目的、有意识的行为，但本质上是客观的、物质的活动。
\begin{itemize}
  \item 首先，构成实践活动的诸要素，即实践的主体、客体和中介都是可感知的客观实在；
  \item 其次，实践的水平、广度、深度和发展过程，都受到客观条件的制约和客观规律的支配；
  \item 最后，实践能够引起客观世界的某种变化，可以把人脑中的观念存在变为现实的存在，给人们提供现实的成果。
\end{itemize}
  \item \textbf{自觉能动性}：人的实践活动是一种有意识、有目的的活动。
  \item \textbf{社会历史性}：实践是社会性的、历史性的活动。
\begin{itemize}
  \item 作为实践主体的人总是处在一定社会关系中；
  \item 实践的社会性决定了它的历史性；
  \item 实践受到一定社会历史条件的制约，并历史地发展着。
\end{itemize}
\end{itemize}
  \item 实践的基本结构：
\begin{itemize}
  \item 人的实践活动是以改造主观世界为目的的客观过程，是实践的主体与客体之间的相互作用，这种相互作用必须借助于一定的手段和工具，即实践的中介。
  \item 实践的主体、客体和中介是实践活动的三项基本要素，三者的有机统一构成实践的基本结构：
\begin{itemize}
  \item 实践主体：在实践活动中，具有一定的主体能力、从事现实社会实践活动的人；
  \item 实践客体：实践活动所指向的对象；
  \item 实践中介：各种形式的工具、手段以及运用、操作这些工具、手段的程序和方法。
\end{itemize}
  \item 实践的主体和客体相互作用的关系，包括实践关系、认识关系和价值关系，其中实践关系是最根本的关系。
  \item 实践的主体、客体和中介是不断变化发展的，因而实践的基本结构也是历史地变化发展的，这种变化主要表现为主体客体化和客体主体化的双向运动。二者作为人类实践活动两个不可分割的方面，互为前提、互为媒介，人类就是通过这种运动形式不断解决着现实世界的矛盾。
\end{itemize}
  \item 实践形式的多样性/实践的三种类型：
\begin{itemize}
  \item 物质生产实践：人类最基本的实践活动，决定着社会的基本性质和面貌；
  \item 社会政治实践：处理各种政治关系的实践，主要指人们的政治活动；
  \item 科学文化实践：创造精神文化产品的实践活动。
\end{itemize}
  \item 实践对认识活动中的决定作用：参见\zhemph{实践是认识的基础}
\end{itemize}
\section{主观能动性和客观规律性的辩证统一}
\begin{itemize}
  \item 正确认识和把握物质和意识的辩证关系，还需要处理好主观能动性和客观规律性的关系。
\begin{itemize}
  \item 一方面，\textbf{尊重客观规律是正确发挥主观能动性的前提}。规律是事物变化发展过程中本身所固有的内在的、本质的、必然的联系。人们只有在认识和掌握客观规律的基础上，才能正确地认识世界，有效地改造世界。
  \item 另一方面，\textbf{只有充分发挥主观能动性，才能正确认识和利用客观规律}。人能够通过自觉活动去认识规律，并按照客观规律去改造世界，以满足自身的需要。
  \item 因此，尊重事物发展的客观规律性与发挥人的主观能动性是辩证统一的，\textbf{实践是客观规律性和主观能动性统一的基础}。
\end{itemize}
  \item 正确发挥人的主观能动性，有以下三个方面的前提和条件：
\begin{itemize}
  \item \textbf{从实际出发}是正确发挥人的主观能动性的前提；
  \item \textbf{实践}是正确发挥人的主观能动性的根本途径；
  \item 正确发挥人的主观能动性，还需要依赖于一定的\textbf{物质条件和物质手段}。
\end{itemize}
\end{itemize}
\section{辩证法的总观点}
世界上的万事万物都处于普遍联系之中，普遍联系引起事物的运动发展。联系和发展的观点是唯物辩证法的总观点，集中体现了唯物辩证法的总特征。
\subsection{事物的普遍联系}
\begin{itemize}
  \item 联系的定义：联系是指事物内部各要素之间和事物之间相互影响、相互制约、相互作用的关系。
  \item 联系的特点：
\begin{itemize}
  \item \textbf{客观性}：联系是事物本身固有的，不是臆想的，要求我们从客观事物本身固有的联系出发去认识事物。
  \item 普遍性：联系的普遍性有三层含义
\begin{itemize}
  \item 其一，任何事物\zhemph{内部的不同部分和不同要素之间}都是相互联系的；
  \item 其二，任何事物都不能孤立存在，都\zhemph{同其他事物}处于一定的联系之中；
  \item 其三，\zhemph{整个世界}是相互联系的统一整体。
\end{itemize}
  \item \textbf{多样性}：分为直接/间接；内部/外部；本质/非本质；必然/偶然联系
  \item \textbf{条件性}：条件是对事物存在和发展发生作用的诸要素的总和，需要辩证看待：
\begin{itemize}
  \item 其一，条件对事物和人的活动具有支持或制约作用；
  \item 其二，条件是可以改变的；
  \item 其三，改变和创造条件并不是任意的，必须尊重事物发展的客观规律。
\end{itemize}
\end{itemize}
\end{itemize}
\subsection{事物的变化发展}
世界上的各种事物不仅是普遍联系的，而且是变化发展的。物质世界处于永恒的运动之中，而物质世界的运动中内在地包含着事物的变化和发展。
\begin{itemize}
  \item 变化泛指事物发生的一切改变，发展则是事物变化中前进的、上升的运动。
  \item 发展的实质是\textbf{新事物的产生和旧事物的灭亡}。新事物是指合乎历史前进方向、具有远大前途的东西，旧事物是指丧失历史必然性，日趋灭亡的东西。
  \item 新事物是不可战胜的：
\begin{itemize}
  \item 新事物与环境：新事物由于具有新的要素、结构和功能，适应已经变化了的环境和条件；
  \item 新事物与旧事物：事物是在旧事物的“母体”中孕育成熟的，对旧事物进行了“扬弃”，并添加了旧事物所不能容纳的新内容，使新事物在本质上优越于旧事物、具有强大生命力。
\end{itemize}
\end{itemize}
\section{对立统一规律与矛盾分析法}
唯物辩证法的基本规律：对立统一规律，量变质变规律，否定之否定规律。对立统一规律是唯物辩证法的实质和核心。
\subsection{对立统一规律与矛盾分析法}
\textbf{对立统一规律是事物发展的根本规律，是唯物辩证法的实质和核心}。对立统一规律提供了人们认识世界和改造世界的根本方法——矛盾分析法。
\begin{mdquote}
矛盾的定义：矛盾是反映事物内部和事物之间对立统一关系的哲学范畴。对立和统一分别体现了矛盾的两种基本属性。矛盾的对立属性又称斗争性，矛盾的统一属性又称同一性。
\end{mdquote}
\begin{itemize}
  \item 矛盾的同一性和斗争性及其在事物发展中的作用：
\begin{itemize}
  \item \textbf{同一性：矛盾着的对立面相互依存、相互贯通的性质和趋势}，有两个方面的含义：
\begin{itemize}
  \item 一是矛盾着的对立面\zhemph{相互依存}，互为存在的前提，并共处于一个统一体中；
  \item 二是矛盾着的对立面相互贯通，在一定条件下可以\zhemph{相互转化}。
\end{itemize}
  \item \textbf{斗争性：矛盾着的对立面相互排斥、相互分离的性质和趋势}，可以分为对抗性矛盾和非对抗性矛盾两种基本形式。
  \item 关系：矛盾的同一性和斗争性相互联结、相辅相成。
\begin{itemize}
  \item 没有斗争性就没有同一性，没有同一性也没有斗争性，斗争性寓于同一性之中，同一性通过斗争性来体现；
  \item 矛盾的同一性是有条件的、相对的，矛盾的斗争性是无条件的、绝对的；
  \item 矛盾的同一性和斗争性相结合，构成了事物的矛盾运动，推动着事物的变化发展。
\end{itemize}
  \item 同一性在事物发展中的作用：
\begin{itemize}
  \item \textbf{同一性是事物存在和发展的前提}；同一性使矛盾双方相互吸取有利于自身的因素，在相互作用中各自得到发展；
  \item \textbf{同一性规定着事物转化的可能和发展的趋势}。
\end{itemize}
  \item 斗争性在事物发展中的作用：
\begin{itemize}
  \item 矛盾双方的斗争促进矛盾双方力量的变化，造成双方力量发展的不平衡，为对立面的转化、事物的质变创造条件；
  \item \textbf{矛盾双方的斗争是一种矛盾统一体向另一种矛盾统一体过渡的决定性力量}。
\end{itemize}
  \item 在事物发展过程中，矛盾的同一性和斗争性相互结合，共同发生作用，但在不同条件下，二者所处的地位会有所不同。在一定的条件下，矛盾的斗争性可能处于主要方面，而在另外的条件下，矛盾的同一性又可能处于主要方面。
  \item 矛盾同一性和斗争性的辩证关系原理，要求我们在观察和处理问题时，必须善于把两者结合起来，在斗争性中把握同一性，在同一性中把握斗争性。和谐是矛盾的一种特殊表现形式，体现着矛盾双方的相互依存、相互存进、共同发展。\textbf{和谐是相对的、有条件的，并不意味着矛盾的绝对同一}。
\end{itemize}
  \item 矛盾的普遍性和特殊性及其相互关系：矛盾既有普遍性，又有特殊性。
\begin{itemize}
  \item 普遍性：矛盾存在于\textbf{一切事物}中，存在于\textbf{一切事物发展过程的始终}；
  \item 特殊性：各个具体事物的矛盾、每一个矛盾的各个方面在发展的不同阶段上各有其特点。
\begin{itemize}
  \item 事物是由多种矛盾构成的。主要矛盾是矛盾体系处于支配地位、对事物发展起决定作用的矛盾。次要矛盾是矛盾体系中处于从属地位、对事物的发展起次要作用的矛盾。
  \item 在每一对矛盾中，处于支配地位、起着主导作用的一方，是矛盾的主要方面；处于被支配地位、不起主导作用的一方，则是矛盾的次要方面。事物的性质是由主要矛盾的主要方面决定的。
  \item 在实际工作中，要坚持“两点论”和“重点论”的统一：
\begin{itemize}
  \item 两点论是指在分析事物的矛盾时，既要看到矛盾双方的对立又要看到统一；既要看到主要矛盾、矛盾的主要方面，又要看到次要矛盾、矛盾的次要方面。
  \item 重点论是指着重把握主要矛盾、矛盾的主要方面，并以此作为解决问题的出发点。
  \item “两点论”和“重点论”的统一要求我们，看问题既要全面地看，又要看主流、大势、发展趋势。
\end{itemize}
\end{itemize}
  \item 矛盾的普遍性和特殊性的辩证统一关系：
\begin{itemize}
  \item 矛盾的普遍性即矛盾的共性，矛盾的特殊性即矛盾的个性。
  \item 矛盾的共性是无条件的、绝对的，矛盾的个性是有条件的、相对的。任何现实存在的事物的矛盾都是共性和个性的有机统一，共性寓于个性之中，没有离开个性的共性，也没有离开共性的个性。
\end{itemize}
\end{itemize}
\end{itemize}
\subsection{量变质变规律}
量变是事物数量的增减和组成要素排列次序的变动，是保持事物的质的相对稳定性的不显著变化，体现了事物发展渐进过程的连续性。质变是事物性质的根本变化，是事物由一种质态向另一种质态的飞跃，体现了事物发展渐进过程和连续性的中断。
\begin{itemize}
  \item 量变和质变的辩证关系：
\begin{itemize}
  \item 量变是质变的必要准备；
  \item 质变是量变的必然结果，并为新的量变开辟道路；
  \item 量变和质变是相互渗透的。
\end{itemize}
\end{itemize}
\subsection{否定之否定规律}
事物的发展是通过其内在矛盾运动以自我否定的方式而实现的。任何事物内部都包含着肯定的方面与否定的方面，由于矛盾双方的相互作用，当否定的方面上升至支配地位时，事物就会由肯定走向对自身的否定，再由否定进一步走向更高阶段的肯定，即否定之否定。否定之否定规律就是要揭示事物自己发展自己的完整过程及本质。
\section{主观辩证法与客观辩证法}
唯物辩证法作为自然、社会和人类思维发展一般规律的科学，是人们认识世界和改造世界的根本方法，是主观辩证法和客观辩证法的有机统一。
\begin{itemize}
  \item 内容：
\begin{itemize}
  \item 客观辩证法：客观事物或客观存在的辩证法，即客观事物以相互作用、相互联系的形式呈现出的各种物质形态的辩证运动和发展规律。
  \item 主观辩证法：人类认识和思维运动的辩证法，即以概念作为思维细胞的辩证思维运动和发展规律。
\end{itemize}
  \item 关系：主观辩证法是客观辩证法在人的思维中的反映，二者在本质上统一，但在表现形式上不同：
\begin{itemize}
  \item 客观辩证法采取\textbf{外部必然性形式}，不以人的意志为转移，是物质世界本身的联系和发展；
  \item 主观辩证法则采取\textbf{观念的、逻辑的形式}，是同人类思维的自觉活动相联系的，是以概念为基础的辩证思维规律，是辩证法的科学体系。
\end{itemize}
\end{itemize}
\section{抽象与具体的思维方法}
抽象与具体是辩证思维的高级形式。在思维活动中，抽象与具体是同分析与综合密切相关的思维方法。这一思维方法是通过从具体到抽象，又从抽象到具体的过程，达到对事物的真理性认识。
\begin{itemize}
  \item 从具体到抽象：
\begin{itemize}
  \item 在认识过程中，有两种完全不同的具体，一种是感性的具体，另一种是思维的具体。
  \item 感性的具体，就是人的感觉器官所得到的知觉表象，是人认识的起点。为了实现从感性的具体到思维的具体的过渡，必须首先否定感性的具体。
  \item 对感性具体的否定就是抽象。抽象是通过分析把整体分解成各个部分，从中抽取出各个必然的、本质的因素，以达到对具体事物的某一本质方面的认识。这是从具体到抽象的过程。
\end{itemize}
  \item 从抽象上升到具体：
\begin{itemize}
  \item 真正达到对具体事物的全面深刻的认识，还必须运用综合的方法，把对事物各方面的本质的认识联系起来，形成对事物整体的统一的认识，使抽象的规定在思维的具体中再现出来。这就是从抽象上升到具体。
\end{itemize}
\end{itemize}
\chapter{第二章\quad{}实践与认识及其发展规律}
\section{实践是认识的基础}
实践是认识的基础，实践在认识活动中起到决定性作用，主要表现在以下四个方面：
\begin{itemize}
  \item 实践是认识的来源：
\begin{itemize}
  \item 认识的内容是在实践活动的基础上产生和发展的；
  \item 离开实践的认识是不可能产生的。
  \item 直接经验与间接经验：一切真知都是从直接经验发源的，但一个人的多数知识还是来自间接经验。
\end{itemize}
  \item 实践是认识发展的动力：
\begin{itemize}
  \item 实践的需要推动认识的产生和发展；
  \item 实践为认识的发展提供了手段和工具；
  \item 实践改造了人的主观世界，锻炼和提高了人的认识能力。
\end{itemize}
  \item 实践是认识的目的：人不是为了认识而认识，其最终目的是为实践服务，指导实践，以满足人们生活和生产的需要。
  \item 实践是检验认识真理性的唯一标准：认识是否具有真理性，只有在实践中才能得到验证。
\end{itemize}
\section{认识的本质及发展规律}
\begin{itemize}
  \item 认识的本质：
\begin{itemize}
  \item 认识的本质是主体在实践基础上对客体的能动反映。这种能动反映不仅具有反映客体内容的反映性特征，而且具有实践所要求的主体能动的、创造性的特征。
  \item 在人的认识活动中，反映特性与能动的创造特性是不可分割的。反映和创造是同一本质的两个不同的方面。创造离不开反映，反映也离不开创造，二者辩证统一。
\end{itemize}
  \item 发展规律：
\begin{itemize}
  \item 人们认识一定事物的过程，是一个从实践到认识，再从认识到实践的过程。
  \item 从实践到认识：
\begin{itemize}
  \item 这个过程主要表现为在实践基础上认识活动由感性认识能动地飞跃到理性认识，这是认识运动的第一次飞跃。这一认知过程可以分为感性认识和理性认识的两个不同阶段：
\begin{itemize}
  \item 感性认识：人们在实践基础上，由感觉器官直接感受到的关于事物的现象、外部联系、各个方面的认识，包括感觉、知觉和表象三种形式；
  \item 理性认识：人们借助抽象思维，在概括整理大量感性材料的基础上，达到关于事物的本质、全体、内部联系和事物自身规律性的认识。
  \item 辩证统一关系：理性认识依赖于感性认识，感性认识有待于发展和深化为理性认识，感性认识和理性认识相互渗透、相互包含，二者辩证统一。
\end{itemize}
\end{itemize}
  \item 从认识到实践：
\begin{itemize}
  \item 认识世界的目的是为了改造世界；
  \item 认识的真理性只有在实践中才能得到检验和发展。
\end{itemize}
\end{itemize}
  \item 实践与认识的辩证运动及其规律：
\begin{itemize}
  \item 实践与认识的辨证运动，是一个由感性认识到理性认识，又由理性认识到实践的飞跃，是实践、认识、再实践、再认识，循环往复以至无穷的辩证发展过程。
  \item 认识是一个反复循环和无限发展的过程。这个过程既是认识在实践基础上沿着科学性方向不断深化发展的过程，也是实践在认识的指导下沿着合理性方向不断深入推进的过程。这个过程既不是封闭式的循环，也不是直线式的发展，往往充满了曲折以至反复，因而是一个波浪式前进和螺旋式上升的过程。
\end{itemize}
\end{itemize}
\section{真理的客观性、绝对性和相对性}
\begin{mdquote}
真理的定义：真理是标志主观和客观相符合的哲学范畴，是对客观事物及其规律的正确反映。
\end{mdquote}
\begin{itemize}
  \item 真理的客观性：
\begin{itemize}
  \item 真理的客观性指真理的内容是对客观事物及其规律的正确反映，真理中包含着不依赖于人和人的意识的客观内容。客观性是真理的本质属性，但真理的形式又是主观的。
  \item 真理是客观的，凡真理都是客观真理，这是真理问题上的唯物论。
\begin{itemize}
  \item 真理的客观性决定了真理的一元性，即同一条件下对于特定的认识客体的真理性认识只有一个。真理是一元的是针对真理的客观内容而言的，真理的主观形式又是多样的。真理是内容上的一元性与形式上的多样性的统一。
  \item 真理的客观性表明，要想发现真理、拥有真理、发展真理并在实践中取得成功，只能采取老老实实的科学态度，尊重真理并按真理办事。
\end{itemize}
\end{itemize}
  \item 真理的绝对性：
\begin{itemize}
  \item 真理的绝对性是指真理主客观统一的确定性和发展的无限性，具有两个方面的含义：
  \item 任何真理都标志着主观与客观相符合，都包含着不依赖于人和人的意识的客观内容，都同谬误有原则的界限，这一点是绝对的、无条件的；
  \item 人类认识按其本性来说，能够正确认识无限发展着的物质世界，认识每前进一步，都是对无限发展着的物质世界的接近，这一点也是绝对的、无条件的。
\end{itemize}
  \item 真理的相对性：
\begin{itemize}
  \item 真理的相对性是指人们在一定条件下对客观事物及其本质和发展规律的正确认识总是有限度的、不完善的，也有两方面的含义：
  \item 从客观世界的整体来看，任何真理都只是对客观世界某一阶段、某一部分的正确认识，人类已经达到的认识的广度总是有限度的$\to$认识有待扩展；
  \item 就特定事物而言，任何真理都只是对客观对象一定方面、一定层次和一定程度的正确认识，认识反映事物的深度是有限度的，或是近似性的$\to$认识有待深化
  \item 因此，任何真理都只能是主观对客观事物近似正确即相对正确的反映。
\end{itemize}
  \item 真理绝对性与相对性的辩证统一关系：
\begin{itemize}
  \item 真理的绝对性和相对性是辩证统一的：
\begin{itemize}
  \item 二者相互依存：任何真理都既是绝对的，又是相对的；
  \item 二者相互包含：绝对之中有相对，相对之中有绝对；
  \item 人类的认识永远处在由真理的相对性走向绝对性、接近绝对性的转化和发展过程中。任何真理性的认识都是由真理的相对性向绝对性转化过程中的一个环节。真理既具有绝对性，又具有相对性，这是真理问题上的辩证法。
  \item 真理的绝对性与相对性根源于人类认识世界能力的无限性与有限性、绝对性与相对性的矛盾。割裂真理的绝对性与相对性的辩证关系，就会走向形而上学的真理观，即绝对主义和相对主义。
\end{itemize}
\end{itemize}
\end{itemize}
\section{真理与谬误}
\begin{itemize}
  \item 真理的定义：真理是标志着主观与客观相符合的哲学范畴，是对客观事物及其规律的正确反映。
  \item 谬误的定义：谬误是同客观事物及其发展规律相违背的认识，是对客观事物及其发展规律的歪曲反映。
  \item 真理与谬误的对立统一关系：真理与谬误是人类认识中的一对永恒矛盾，他们之间既对立又统一：
\begin{itemize}
  \item 真理与谬误相互对立。在确定的对象和范围内，其对立是绝对的，存在着原则界限。
  \item 真理与谬误的对立又是相对的，在一定条件下能够相互转化：
\begin{itemize}
  \item 首先，真理在一定条件下会转化为谬误：
\begin{itemize}
  \item 一方面，真理是具体的。任何真理都是在一定范围内、一定条件下才能够成立，超出这个范围，失去特定条件，就会变成谬误；
  \item 另一方面，真理又是全面的。把全面的真理性认识组成的科学体系中的某个原理孤立地抽取出来，切断同其他原理的联系，也会使其丧失自己的真理性而变为谬误。
\end{itemize}
  \item 其次，谬误在一定条件下能够向真理转化。
  \item 最后，在批判谬误中发展真理，是谬误向真理转化的另一种形式。
\end{itemize}
  \item 真理和谬误的对立统一关系表明，真理总是同谬误相比较而存在、相斗争而发展的。
\end{itemize}
\end{itemize}
\section{实践是检验真理的唯一标准}
实践之所以能够成为检验真理的唯一标准，是由真理的本性和实践的特点决定的。
\begin{itemize}
  \item 第一，从真理的本性看，真理是人们对客观事物及其发展规律的正确反映，它的本性在于主管和客观相符合。检验真理就是检验人的主观认识同客观实际是否相符合以及符合的程度。
\begin{itemize}
  \item 检验真理的标准，既不能是主观认识本身，也不能是客观事物。只有那种能够把主观认识与客观事物联系和沟通起来，从而使人们能够把二者加以比较和对照的东西，才能充当检验真理的标准。具有这种特性的东西，只能是作为主客观联系的桥梁的社会实践。
\end{itemize}
  \item 第二，从实践的特点来看，实践具有直接现实性。实践的直接现实性是它的客观实在性的具体表现。
\begin{itemize}
  \item 如果实践的结果与实践之前的认识相符合，那么之前的认识就得到了证实，成为了真理性的证实。
  \item 实践的直接现实性的品格，是实践能够成为检验真理唯一标准的主要根据。
\end{itemize}
  \item 在实践检验真理的过程中，逻辑证明可以起到重要的补充作用。但是，逻辑证明只能回答前提与结论的关系是否符合逻辑的问题，而不能回答结论是否符合客观实际的问题。已被逻辑证明了的东西，还必须经过实践的检验，并最终服从实践检验的结果。
\end{itemize}
\section{必然与自由}
\begin{itemize}
  \item 自由的定义：自由是表示人的活动状态的范畴，是指人在活动中通过认识和利用必然所表现出的一种自觉自主的状态。
  \item 必然的定义：必然性即规律性，指的是不依赖于人的意识而存在的自然和社会发展所固有的客观规律。
  \item 必然与自由的辨证规律：
\begin{itemize}
  \item 人不能摆脱必然性的制约，只有在认识必然性的基础上才有自由的活动，这就是人的自由限度，也是自由和必然的辩证规律。
  \item 自由是有条件的。不具备认识条件和实践条件，就得不到自由。
\begin{itemize}
  \item 认识条件：要有对客观事物运动发展规律性、必然性的正确认识；
  \item 实践条件：能够将获得的规律性认识运用于指导实践，实现改造世界的目的。
\end{itemize}
  \item 认识必然和争取自由，是人类认识世界和改造世界的根本目标，是一个历史性的过程。必然与自由的关系贯穿于人类存在和发展的始终，并成为人类存在和发展的永恒矛盾，因此也是人类存在和发展的永恒动力。
\end{itemize}
\end{itemize}
\chapter{第三章\quad{}人类社会及其发展规律}
\section{社会存在和社会意识的辨证关系}
\begin{itemize}
  \item 社会存在：社会物质生活条件，是社会生活的物质方面，主要包括自然地理环境、人口因素和物质生产方式。
  \item 社会意识：社会存在的反映，是社会生活的精神方面。
\begin{itemize}
  \item 按照不同主体，可分为个体意识和群体意识；
  \item 按照不同层次，可分为低层次的社会心理和高层次的社会意识形式。其中，社会意识形式存在意识形态和非意识形态之分。社会意识形式以社会心理为基础，并对社会心理起指导和影响作用。
\end{itemize}
  \item 社会存在和社会意识的关系是辩证统一的。\textbf{社会存在决定社会意识，社会意识是社会存在的反映，并反作用于社会存在}。
\begin{itemize}
  \item 社会存在决定社会意识：
\begin{itemize}
  \item 社会存在是社会意识内容的客观来源，社会意识是社会物质生活过程及其条件的主观反映；社会意识根源于社会存在，是对以实践为基础的不断变化的现实世界的反映；社会意识是人们进行社会物质交往的产物；
  \item 社会意识是具体的、历史的。随着社会存在的发展，社会意识也相应地或早或迟地变化和发展，但不管怎样变化和发展，其根源总是深藏于经济事实当中。
\end{itemize}
  \item 社会意识反映并反作用于社会存在：
\begin{itemize}
  \item 社会意识以理论、观念、心理等形式反映社会存在，这是社会意识对社会存在的依赖性。
  \item 但社会意识并非消极被动地受制于社会存在，它既依赖于社会存在，又有其相对独立性，具有自身特有的发展形式和发展规律。表现为：
\begin{itemize}
  \item 社会意识与社会存在发展具有不完全同步性和不平衡性；
  \item 社会意识内部各种形式之间存在相互影响且各自具有历史继承性；
  \item 社会意识对社会存在具有能动的反作用，这是社会意识相对独立性的突出表现。
\end{itemize}
  \item 社会意识的能动作用是通过指导人们的实践活动实现的。
\end{itemize}
\end{itemize}
\end{itemize}
\section{社会基本矛盾及其运动的规律}
\begin{itemize}
  \item 生产力与生产关系、经济基础与上层建筑之间的矛盾，是人类社会基本矛盾。
  \item 生产力与生产关系矛盾运动的规律、经济基础与上层建筑矛盾运动的规律，是人类社会发展的基本规律。
\end{itemize}
\subsection{生产力与生产关系的矛盾运动及其规律}
\begin{itemize}
  \item 什么是生产力：
\begin{itemize}
  \item 生产力是人类在生产实践中形成的改造和影响自然以使其适合社会需要的物质力量，具有客观现实性和社会历史性。
  \item 生产力是结构复杂的系统，其基本要素包括：
\begin{itemize}
  \item 劳动资料，又称劳动手段
  \item 劳动对象
  \item 劳动者
\end{itemize}
  \item 科学技术是生产力中的重要因素，是先进生产力的集中体现和主要标志，是第一生产力。
\end{itemize}
  \item 什么是生产关系：
\begin{itemize}
  \item 生产关系指的是人们在物质生产过程中形成的不以人的意识为转移的经济关系，是社会关系中最基本的关系。
  \item 生产关系包括生产资料所有制关系、生产中人与人的关系和产品的分配关系。其中，\textbf{生产资料所有制关系是最基本的、最有决定意义的方面}。
  \item 依据生产资料所有制关系的性质，生产关系区分为两种基本类型：以生产资料的公有制为基础的生产关系，以及以生产资料私有制为基础的生产关系。随着生产力的发展，前者必将取代后者。
\end{itemize}
  \item 生产关系一定要适合生产力状况的规律：
\begin{itemize}
  \item 生产力和生产关系是社会生产不可分割的两个方面，在社会生产中，生产力是生产的物质内容，生产关系是生产的社会形式，二者的有机统一构成社会的生产方式。
  \item 生产力与生产关系的相互关系是：\textbf{生产力决定生产关系，而生产关系又反作用于生产力}。
\begin{itemize}
  \item 第一，生产力决定生产关系。在二者的关系中，生产力是居支配地位、起决定作用的方面。
\begin{itemize}
  \item 生产力状况决定生产关系的性质；
  \item 生产力的发展决定生产关系的变化。
\end{itemize}
  \item 第二，生产关系对生产力具有能动的反作用。
\begin{itemize}
  \item 当生产关系适合生产力发展的客观要求时，对生产力的发展起推动作用；
  \item 当生产关系不适合生产力发展的客观要求时，就会阻碍生产力的发展。
\end{itemize}
\end{itemize}
  \item 生产力与生产关系的相互作用是一个过程，表现为二者的矛盾运动。这种矛盾运动中内在的、本质的、必然的联系，就是生产关系一定要适合生产力状况的规律。
\end{itemize}
\end{itemize}
\subsection{经济基础与上层建筑的矛盾及其规律}
\begin{itemize}
  \item 什么是经济基础：
\begin{itemize}
  \item 定义：经济基础是指由社会一定发展阶段的生产力所决定的生产关系的总和。
  \item 内涵：
\begin{itemize}
  \item 社会的一定发展阶段往往存在多种生产关系，但决定一个社会性质的是其中占支配地位的生产关系；
  \item 经济基础与经济体制具有内在联系。
\end{itemize}
\end{itemize}
  \item 什么是上层建筑：
\begin{itemize}
  \item 上层建筑是建立在一定经济基础之上的意识形态以及与之相适应的制度、组织和设施。
  \item 构成：
\begin{itemize}
  \item 意识形态（观念上层建筑）
  \item 政治法律制度及设施和政治组织（政治上层建筑）
  \par 政治上层建筑居于主导地位，国家政权是政治上层建筑的核心。
\end{itemize}
  \item 上层建筑一定要适合经济基础状况的规律：
\begin{itemize}
  \item 经济基础和上层建筑是辩证统一的。\textbf{经济基础决定上层建筑，上层建筑反作用于经济基础}，二者相互影响、相互作用：
\begin{itemize}
  \item 首先，经济基础决定上层建筑：
\begin{itemize}
  \item 经济基础是上层建筑赖以产生、存在和发展的物质基础；
  \item 上层建筑是经济基础得以确立其统治地位并获得巩固和发展不可缺少的政治、思想条件。
\end{itemize}
  \item 其次，上层建筑对经济基础具有反作用：
\begin{itemize}
  \item 表现：上层建筑为自己的经济基础的形成和巩固服务，确立或维护其在社会中的统治地位。
  \item 后果：当上层建筑为适合生产力发展要求的经济基础服务时，就成为推动社会发展的进步力量；当它为落后的经济基础服务时，就成为阻碍社会发展的消极力量。
\end{itemize}
  \item 再次，经济基础和上层建筑之间的相互作用构成二者的矛盾运动。
\begin{itemize}
  \item 其一，在同一性质的经济基础与上层建筑的关系中，上层建筑的不完善部分、没有反映经济基础要求的部分都会同经济基础发生矛盾；
  \item 其二，在不同性质的经济基础与上层建筑的关系中，矛盾更为复杂，主要表现在：占统治地位的经济基础同旧上层建筑的残余、未来上层建筑的萌芽之间的矛盾，新旧上层建筑之间、新旧经济基础之间的矛盾等；
  \item 其三，当一种社会形态处于上升发展阶段时，上层建筑对于经济基础一般是适应的；当一种社会形态处于没落时期时，上层建筑同经济基础变革的客观要求是不相适应的，其矛盾则变为对抗性的、全局性的矛盾。
  \item 最后，经济基础和上层建筑之间的内在联系构成了上层建筑一定要适合经济基础状况的规律。
\begin{itemize}
  \item 经济基础状况决定上层建筑的发展方向，决定上层建筑相应的调整或变革；
  \item 上层建筑的反作用也必须取决于和服从于经济基础的性质和客观要求。
\end{itemize}
\end{itemize}
\end{itemize}
\end{itemize}
\end{itemize}
\end{itemize}
\section{社会形态}
\begin{itemize}
  \item 社会形态的内涵：
\begin{itemize}
  \item 社会形态是关于社会运动的具体形式、发展阶段和不同质态的范畴，是同生产力发展一定阶段相适应的经济基础与上层建筑的统一体。
  \item 社会形态包括社会的\zhemph{经济形态、政治形态和意识形态}，是三者的具体的、历史的统一。
  \item 一定的社会形态总要以一定的社会制度形式呈现出来，社会制度能够集中体现社会形态的性质，所以人们常用社会制度指代社会形态。
  \item 人类社会是不断发展的，社会的根本性变革和进步就是通过社会形态的更替实现的。
\end{itemize}
  \item 社会形态更替的统一性和多样性：
\begin{itemize}
  \item 依据生产关系的不同性质，社会历史可划分为五种社会形态：原始社会、奴隶社会、封建社会、资本主义社会和共产主义社会。这五种社会形态的依次更替，是社会历史运动的一般过程和一般规律，表现了社会形态更替的统一性。
  \item 但是就某一国家或民族的社会发展历程而言，情况就复杂了。不同国家社会形态更替的过程不同；即便是同一种意识形态，在不同国家也会呈现出不同特点。所有这些都体现了社会形态更替形式的多样性。
\end{itemize}
  \item 社会形态更替的必然性和人们的历史选择性：
\begin{itemize}
  \item 社会形态更替的统一性和多样性，根源于社会发展的客观必然性与人们的历史选择性相统一的过程。
  \item 社会形态更替的客观必然性主要是指社会形态依次更替的过程和规律是客观的，其发展的基本趋势是确定不移的。生产力与生产关系矛盾运动的规律性，从根本上规定了社会形态更替的客观必然性。
  \item 社会形态更替的规律也是人们自己的社会行动的规律。规律的客观性并不否认人们历史活动的能动性，并不排斥人们在遵循社会发展规律的基础上，对于某种社会形态的历史选择性：
\begin{itemize}
  \item 社会发展的客观必然性造成了一定历史阶段社会发展的基本趋势，为人们的历史选择提供了基础、范围和可能性空间；
  \item 社会形态更替的过程也是一个\textbf{主观能动性和客观规律性相统一}的过程；
  \item 人们的历史选择性归根结底是人民群众的选择性。人们对于社会形态的历史选择最终取决于人民群众的根本利益、根本意愿以及对社会发展规律的把握和顺应程度。
\end{itemize}
  \item 社会形态的更替的前进性与曲折性：
\begin{itemize}
  \item 社会形态的更替还表现为历史的前进性与曲折性、顺序性与跨越性的统一。
  \item 社会形态更替的前进性、顺序性主要是指五种社会形态依次演进的基本趋势，其历史过程是一个“扬弃”的过程，但它并不否认历史发展的曲折性和跨越性。
  \item 在社会进步发展的过程中，有时会出现社会形态更替的反复甚至倒退现象。
\end{itemize}
\end{itemize}
\end{itemize}
\section{社会发展的根本动力}
\textbf{社会历史发展的根本动力是社会基本矛盾}。社会基本矛盾是指贯穿社会发展过程始终、规定社会发展过程的基本性质和基本趋势，并对社会历史发展起根本推动作用的矛盾。\textbf{生产力和生产关系、经济基础和上层建筑的矛盾是社会基本矛盾}。
\begin{itemize}
  \item 社会基本矛盾在历史发展中的作用主要表现在：
\begin{itemize}
  \item 首先，\textbf{生产力是社会基本矛盾运动中最基本的动力因素}，是人类社会发展和进步的最终决定力量；生产力是社会进步的根本内容，是衡量社会进步的根本尺度。
  \item 其次，社会基本矛盾特别是生产力和生产关系的矛盾，\textbf{决定着社会中其他矛盾的存在和发展}；经济基础和上层建筑的矛盾也会影响和制约生产力和生产关系的矛盾。
  \item 最后，\textbf{社会基本矛盾具有不同的表现形式和解决方式，并从根本上影响和促进社会形态的变化和发展}。
\end{itemize}
\end{itemize}
\section{人民群众是历史的创造者}
\subsection{唯物史观认为人民群众是历史的创造者}
唯物史观与唯心史观的对立，在历史创造者的问题上表现为群众史观与英雄史观的对立。英雄史观从\textbf{社会意识决定社会存在}的基本前提出发，否认\textbf{物质资料的生产方式是社会发展的决定力量}，抹杀人民群众的历史作用，宣扬少数英雄人物创造历史。相反，群众史观认为历史的创造者不是个别英雄，而是人民群众。
\begin{itemize}
  \item 首先，唯物史观立足于\textbf{现实的人及其本质}来把握历史的创造者；
  \item 其次，唯物史观立足于\textbf{整体的社会历史过程}来探究谁是历史的创造者；
  \item 再次，唯物史观从\textbf{社会历史发展的必然性}入手考察和说明谁是历史的创造者；
  \item 最后，唯物史观从\textbf{人与历史关系的不同层次}上考察谁是历史的创造者。它不是对历史表象的经验描述，而是对历史本质的逻辑把握。
\end{itemize}
\subsection{人民群众在创造历史过程中的决定作用}
\begin{mdquote}
\begin{itemize}
  \item 从质上看，人民群众是指一切\zhemph{对社会历史发展起推动作用的人}；
  \item 从量上看，人民群众是指社会人口的绝大多数。
  \item 人民群众是一个历史范畴。在当代中国，全体社会主义劳动者、社会主义事业的建设者、拥护社会主义的爱国者、拥护祖国统一和致力于中华民族伟大复兴的爱国者都属于人民群众的范畴。
\end{itemize}
\end{mdquote}
在社会历史发展过程中，人民群众起着决定性的作用：
\begin{itemize}
  \item 人民群众是社会历史实践的主体，在创造历史中起决定性的作用；
  \item 人民群众是社会物质财富的创造者；
  \item 人民群众是社会精神财富的创造者；
  \item 人民群众是社会变革的决定力量。
  \item 人民群众的历史活动受到一定社会历史条件的制约。
\begin{itemize}
  \item 经济条件对于人民群众创造历史的活动有着首要的、决定性的影响。
\end{itemize}
\end{itemize}
\subsection{无产阶级政党的群众路线}
\begin{itemize}
  \item 唯物史观关于人民群众是历史创造者的原理，要求我们坚持马克思主义群众观点，贯彻党的群众路线：
\begin{itemize}
  \item 群众观点：坚信人民群众自己解放自己的观点，一切向人民群众负责的观点，虚心向群众学习的观点。
  \item 群众路线：群众路线是群众观点的具体应用，即一切为了群众、一切依靠群众，从群众中来，到群众中去。实质在于充分相信群众，坚决依靠群众，密切联系群众，全心全意为人民群众服务。
\end{itemize}
\end{itemize}
\subsection{个人在社会历史中的作用}
任何历史人物的出现都体现了必然性与偶然性的统一。
\begin{itemize}
  \item 时势造英雄，杰出人物的出现具有必然性。
  \item 杰出人物对社会进程的产生影响仅仅是历史进程中的偶然现象，只能成为社会发展的个别原因。它们凭借自己的才能，虽然也能使具体历史事件的外貌或某些后果发生变化，但终究不能改变历史发展的基本方向。
  \item 如果看不到历史人物活动的社会制约性，割裂偶然与必然的关系，就势必会夸大个人的作用，进而否定或歪曲历史发展的规律。
\end{itemize}
\chapter{第四章\quad{}资本主义的本质及规律}
\section{商品经济产生的条件}
\begin{itemize}
  \item 商品经济是以交换为目的而进行生产的经济形式，是一定社会历史条件的产物。
  \item 商品经济得以产生的社会历史条件有两个：
\begin{itemize}
  \item \textbf{存在社会分工}；
  \item \textbf{生产资料和劳动产品属于不同的所有者}。
\end{itemize}
\end{itemize}
\section{商品二因素与劳动二重性}
\begin{itemize}
  \item 商品二因素：商品是用来交换、能满足人的某种需要的劳动产品，具有使用价值和价值两个因素或两种属性，是使用价值和价值的矛盾统一体。
\begin{itemize}
  \item 使用价值：商品能满足人的某种需要的有用性，反映的是人与自然之间的物质关系，是商品的自然属性，是一切劳动产品所共有的属性，离开了它商品就不复存在。
  \item 价值：凝结在商品中的无差别的人类劳动，即人的脑力和体力的耗费，价值是商品所特有的社会属性。
  \item 使用价值和价值之间对立统一的关系：
\begin{itemize}
  \item 对立性：商品的使用价值和价值是相互排斥的，二者不可兼得。要获得商品的价值就必须放弃商品的使用价值；要得到商品的使用价值，就不能得到商品的价值。
  \item 统一性：作为商品必须同时具有使用价值和价值两个因素。使用价值是价值的物质承担者，价值寓于使用价值之中。
\end{itemize}
\end{itemize}
  \item 劳动二重性：
\begin{itemize}
  \item 生产商品的劳动具有二重性，即具体劳动和抽象劳动。
\begin{itemize}
  \item 具体劳动：生产一定使用价值的具体形式的劳动。
  \item 抽象劳动：撇开一切具体形式的、无差别的人类劳动，即人的脑力和体力的耗费。
\end{itemize}
  \item 生产商品的具体劳动创造商品的使用价值，抽象劳动形成商品的价值。任何一种劳动，一方面是特殊的具体劳动，另一方面又是一般的抽象劳动，这就是劳动的二重性。正是劳动的二重性决定了商品的二因素。
  \item 具体劳动和抽象劳动也是对立统一的关系：
\begin{itemize}
  \item 一方面，具体劳动和抽象劳动不是各自独立存在的两种劳动或两次劳动，它们在时间和空间上是统一的，是商品生产者的同一劳动过程不可分割的两个方面；
  \item 另一方面，具体劳动与抽象劳动又分别反映劳动的不同属性，具体劳动所反映的是人与自然的关系，是劳动的自然属性，而抽象劳动所反映的是商品生产者的社会关系，是劳动的社会属性。
\end{itemize}
\end{itemize}
\end{itemize}
\section{价值量及其决定；社会必要劳动时间；劳动生产率}
\begin{itemize}
  \item 商品的价值与价值量：
\begin{itemize}
  \item 商品的价值是凝结在商品中的无差别的一般人类劳动，价值量是由劳动者生产商品所耗费的劳动量决定的，而劳动量则按照劳动时间来计量。
  \item 决定商品价值量的不是生产商品的个别劳动时间，而是社会必要劳动时间。
\end{itemize}
  \item 社会必要劳动时间：
\begin{itemize}
  \item 社会必要劳动时间是在现有的社会正常的生产条件下，在社会平均的劳动熟练程度和劳动强度下制造某种使用价值所需要的劳动时间，即形成价值量的劳动力和生产条件都必须具有正常的性质。
\end{itemize}
  \item 劳动生产率：
\begin{itemize}
  \item 劳动生产率指的是劳动者生产使用价值的效率。它的高低可以用单位劳动时间内生产的商品数量来测量，也可以用单位商品中所耗费的劳动时间来测量。
  \item 生产商品所需要的社会必要劳动时间随着劳动生产率的变化而变化。劳动生产率水平越高，单位时间内生产的商品数量越多，生产每件商品所需要的社会必要劳动时间就越少，单位商品的价值量就越小，反之就越大。商品的价值量与生产商品所耗费的社会必要劳动时间成正比，与劳动生产率成反比。
  \item 影响劳动生产率的因素很多，主要有劳动者的平均熟练程度、科学技术的发展水平及其在生产中的应用程度、生产过程的社会结合、生产资料的规模和效能以及自然条件等。
\end{itemize}
\end{itemize}
\section{货币起源；货币的本质和职能}
\begin{itemize}
  \item 货币起源：从历史上看，商品价值形式的发展经历了四个阶段，即简单的或偶然的价值形式、总和的或扩大的价值形式、一般价值形式、以及货币形式。
  \item 货币的本质：货币是在长期交换过程中形成的固定充当一般等价物的商品，是商品经济内在矛盾发展的产物，货币的本质体现一种社会关系。
  \item 货币的职能：货币具有价值尺度、流通手段、贮藏手段、支付手段和世界货币等职能，其中价值尺度和流通手段是最基本的职能。
\end{itemize}
\section{价值规律及其作用}
\begin{itemize}
  \item 定义：价值规律是商品生产和商品交换的基本规律。价值规律贯穿于商品经济的全部过程，它既支配商品生产，又支配商品流通。
  \item 主要内容与客观要求：商品的价值量由生产商品的社会必要劳动时间决定，商品交换以价值量为基础，按照等价交换的原则进行。
  \item 表现形式：商品的价格围绕商品的价值自发波动。由于供求关系变化的影响，商品价格总是时而高于价值，时而低于价值，不可避免地围绕价值这个中心上下波动。从较长时间来看，价格高于价值的部分和价格低于价值的部分能够相抵，商品的平均价格和价值是相一致的。
  \item 在市场配置资源过程中的作用：
\begin{itemize}
  \item 自发地调节生产资料和劳动力在社会各生产部门之间的分配比例；自发地刺激社会生产力的发展；
  \item 自发地调节社会收入分配。
  \item 消极后果：导致资源浪费；阻碍技术进步；导致收入两极分化。
\end{itemize}
\end{itemize}
\section{商品经济的基本矛盾}
\begin{mdquote}
\textbf{私人劳动和社会劳动的矛盾是商品经济的基本矛盾。}商品的使用价值和价值的矛盾、生产商品的具体劳动和抽象劳动的矛盾，根源于私人劳动和社会劳动的矛盾。
\end{mdquote}
\begin{itemize}
  \item 在以\textbf{私有制为基础的商品经济}中，\textbf{商品生产者的劳动既是具有社会性质的社会劳动，又是具有私人性质的私人劳动}：
\begin{itemize}
  \item 商品生产者的劳动的\textbf{社会性质是由社会分工决定的}。每个商品生产者的劳动都是社会总劳动的一部分，是具有社会性质的社会劳动。
  \item 商品生产者的劳动的\textbf{私人性质是由生产资料私有制决定的}。商品生产者的劳动是按照自己的利益和要求进行的，是具有私人性质的私人劳动。
\end{itemize}
  \item 私人劳动和社会劳动的矛盾构成私有制商品经济的基本矛盾，这一矛盾贯穿商品经济发展过程的始终，决定着商品经济的各种内在矛盾及其发展趋势：
\begin{itemize}
  \item 首先，私人劳动和社会劳动的矛盾决定着\textbf{商品经济的本质及发展过程}；
\begin{itemize}
  \item 商品经济是以交换为目的的经济形式，交换体现了商品经济的本质；
  \item 交换是解决私人劳动和社会劳动之间矛盾的唯一途径。
\end{itemize}
  \item 其次，私人劳动和社会劳动的矛盾是商品经济其他一切矛盾的基础；
\begin{itemize}
  \item 在商品交换过程中，只有将具体劳动还原为抽象劳动，才能进行量的比较。二具体劳动能否还原为抽象劳动，在根本上取决于私人劳动和社会劳动能否实现统一。
  \item 商品生产者的私人劳动生产的产品如果与社会的需求不相适应，这种私人劳动就不会被承认为社会劳动，它作为具体劳动的有用性质也就不会为社会所承认，因而不能还原为抽象劳动，这意味着商品的价值无法实现，商品的使用价值和价值之间的矛盾没有得到解决。
\end{itemize}
  \item 最后，私人劳动和社会劳动的矛盾决定着商品生产者的命运。
\begin{itemize}
  \item 商品的售卖过程是私人劳动转化为社会劳动的过程。
\end{itemize}
\end{itemize}
  \item 私有制商品经济条件下，私人劳动和社会劳动之间的矛盾是商品经济的基本矛盾。在资本主义制度下，这种矛盾进一步发展成资本主义的基本矛盾，即\textbf{生产社会化和生产资料资本主义私人占有之间的矛盾}，正是这一矛盾的不断运动，使资本主义制度最终被社会主义制度所替代具有了客观必然性。
\end{itemize}
\section{商品拜物教}
\begin{mdquote}
私有制商品经济条件下\textbf{私人劳动与社会劳动之间的矛盾}通过商品的运动、 价值的运动、 货币的运动决定商品生产者的命运， 这使商品生产者认为商品、 价值乃至货币似乎是物的\textbf{自然属性}， 而这种所谓的自然属性又似乎具有一种\textbf{超自然的神秘性}， 商品生产者不能自己掌握自己的命运， 而是听凭商品、 价值、 货币运动的摆布， \textbf{人与人之间一定的社会关系在人们面前采取了物与物的关系的虚幻形式}， 马克思称之为 “商品拜物教＂ 。
\end{mdquote}
私有制商品经济条件下商品世界的拜物教性质的产生具有必然性。
\begin{itemize}
  \item 其一，\textbf{劳动产品只有采取商品的形式才能进行交换}， 人类劳动的等同性只有采取同质的价值形式才能在交换中体现出来。 
  \item 其二， \textbf{劳动量只有采取价值量这一物的形式才能进行计算和比较}。 
  \item 其三， \textbf{生产者的劳动关系的社会性质只有采取商品之间即物与物之间相交换的形式才能间接地表现出来}， 这就使人与人之间一定的社会关系被物与物的关系所掩盖， 具有了拜物教性质。 商品世界的拜物教性质或者生产商品的劳动的社会性质所具有的物的外观掩盖了商品经济关系的本质， 妨碍人们透过物的外表认识商品、 价值以及货币的实质。
\end{itemize}
\section{资本主义生产关系的产生和资本主义生产方式的形成}
\subsection{资本主义生产关系的产生（雇佣关系）}
\begin{itemize}
  \item 背景：随着社会生产力的提高特别是封建社会后期商品经济的迅速发展，封建生产关系越来越成为生产力发展的障碍。生产力与生产关系的矛盾引起了农村自然经济和城市行会组织的瓦解，导致城乡资本主义生产关系的产生。
  \item 途径：资本主义萌芽于14世纪末期15世纪初地中海沿岸的一些城市，资本主义的产生途径有两个：一是从小商品经济分化出来，二是从商人和高利贷者分化出来。
\begin{itemize}
  \item 随着商品经济的发展，小生产者之间展开激烈的竞争并逐渐分化，一部分条件较好的作坊主成为最早的工业资本家，而大多数作坊主在竞争中衰落直至破产，最终同帮工和学徒沦为雇佣工人。手工作坊中的师徒关系逐渐转变为雇佣关系，成为资本主义生产关系萌芽的形式之一。
  \item 随着商人积累的财富不断增加，一些大商人成了包买商，割断了小生产者与销售市场和原料市场的联系，逐渐控制了商品生产者。商人和高利贷者趁生产者困难之机，贷给他们所需要的资金、原料和生产工具，生产者一旦无力还债，就只能交出自己的作坊抵债。于是，作坊主丧失了独立的生产者身份，连同其帮工和学徒成了商人或高利贷者的雇佣工人，商人或高利贷者则成为工业资本家，这成为资本主义生产关系萌芽的另一种形式。
\end{itemize}
\end{itemize}
\subsection{资本主义生产方式的形成（资本原始积累）}
\begin{itemize}
  \item 定义：资本原始积累就是以暴力手段使生产者与生产资料相分离，资本迅速集中于少数人手中，资本主义得以迅速发展的历史过程。
  \item 途径：资本积累主要是通过两个途径进行的：
\begin{itemize}
  \item 用暴力手段剥夺农民的土地，是资本原始积累过程的基础，在英国表现得最为典型。资本家和封建贵族通过各种手段将农民私有土地改为养羊的牧场，农民变成一无所有的流浪者，不得不到资本家的工厂出卖劳动力。
  \item 用暴力手段掠夺货币财富的一个重要表现就是利用国家政权的力量进行残酷的殖民掠夺，新兴资产阶级武力征服海外殖民地、掠夺大量财富，大大加速了货币资本的积累。这一切都大大促进了资本主义的发展，缩短了封建生产方式转变为资本主义生产方式的历史进程。
\end{itemize}
\end{itemize}
\section{劳动力成为商品与货币转化为资本}
\begin{mdquote}
劳动力是指人的劳动能力，是人的脑力和体力的总和。劳动力的使用即劳动。
\end{mdquote}
\begin{itemize}
  \item 劳动力成为商品的基本条件：
\begin{itemize}
  \item 其一，劳动力是自由人，\textbf{能把自己的劳动力当作资本来支配}；
  \item 其二，\textbf{劳动者没有任何生产资料，没有生活资料来源，因而不得不依靠出卖劳动力为生}。
\end{itemize}
  \item 劳动力商品的特点：
\begin{itemize}
  \item 和任何商品一样，劳动力商品也有它的价值和使用价值。但是劳动力是特殊的商品，它的价值和使用价值具有不同于普通商品的特点。
  \item 劳动力的价值是由生产、发展、维持和延续劳动力所必需的生活必需品的价值决定的，劳动力价值的最低界限是由生活上不可缺少的生活资料的价值决定的。一旦劳动力价值降低到这个界限以下，劳动力就只能在萎缩的状态下维持。
  \item （劳动力成为商品是货币成为资本的前提条件）劳动力商品在使用价值上有一个很突出的特点，就是它的\textbf{使用价值是价值的源泉}，在消费过程中能够创造新的价值，而且这个新的价值比劳动力本身的价值更大。正是由于这一特点，货币所有者购买到劳动力以后，在消费过程中，不仅能够收回他在购买这种特殊商品时支付的价值，还能得到一个增殖的价值即剩余价值($m$)。而\textbf{一旦货币购买的劳动力带来剩余价值，货币也就变成了资本}。
\end{itemize}
\end{itemize}
\section{资本主义所有制}
\begin{itemize}
  \item 含义：
\begin{itemize}
  \item 资本主义所有制是生产资料归资本家所有的一种私有制形式。
  \item 在资本主义所有制条件下，资本家不但拥有生产资料的所有权，而且拥有对雇佣劳动者的支配权，并凭借这种所有权和支配权实现对全部劳动产品的占有和支配。
\end{itemize}
  \item 本质：
\begin{itemize}
  \item 资本家凭借对生产资料的占有，在\zhemph{等价交换原则的掩盖下}，雇佣工人从事劳动，\textbf{无偿占有雇佣工人创造的剩余价值}，资本与雇佣劳动的关系由此具有了剥削与被剥削的对抗性质。
  \item 因此，资本主义所有制是雇佣劳动赖以存在的基础，是资本与雇佣劳动之间剥削与被剥削关系的体现。这就是资本主义所有制的本质。
\end{itemize}
\end{itemize}
\section{生产剩余价值是资本主义生产方式的绝对规律}
\subsection{剩余价值的生产过程}
\begin{itemize}
  \item 剩余价值是在资本主义生产过程中生产出来的。资本主义的生产过程具有二重性，一方面是生产物质资料的劳动过程；另一方面是生产剩余价值的过程，即价值增殖过程。资本主义生产过程是劳动过程和价值增殖过程的统一。
  \item 价值增殖过程是剩余价值的生产过程，这是资本主义生产过程的主要方面。所谓价值增殖过程，是超过劳动力价值的补偿这个一定点而延长了的价值形成过程。
\begin{itemize}
  \item 如果劳动者创造的价值刚好补偿资本家所预付的劳动力价值，那就是单纯的价值形成过程；
  \item 如果价值形成过程超过了这个一定点，就变成了价值增殖过程。
\end{itemize}
  \item 在价值增殖过程中，雇佣工人的劳动分为两部分，一部分是必要劳动，用于再生产劳动力的价值；另一部分是剩余劳动，用于无偿地为资本家生产剩余价值。因此，剩余价值是雇佣工人所创造的并被资本家无偿占有的超过劳动力价值的那部分价值，它是雇佣工人剩余劳动的凝结，体现了资本家与雇佣工人之间剥削和被剥削的关系。
\end{itemize}
\subsection{资本的本质}
\begin{itemize}
  \item 资本是能够带来剩余价值的价值。剩余价值是由雇佣工人的剩余劳动创造的。
  \item 在资本主义社会，资本总是通过各种物品表现出来，但\zhemph{资本不是物}，而是一定的、社会的、属于一定历史社会形态的\zhemph{生产关系}，后者体现在一个物上，并赋予这个物以独特的社会性质。
\end{itemize}
\subsection{不变资本与可变资本}
资本在资本主义生产过程中采取生产资料和劳动力两种形态，根据这两部分资本在\textbf{剩余价值生产中所起的不同作用}，可将资本区分为不变资本和可变资本。
\begin{itemize}
  \item 不变资本($c$)：不变资本是以\textbf{生产资料}形态存在的资本。生产资料的价值通过工人的具体劳动被转移到新产品中，其转移的价值量不会大于它原有的价值量；
  \item 可变成本($v$)：可变成本是用来\textbf{购买劳动力}的那部分成本。可变资本的价值在生产过程中不是被转移到新产品上去，而是由工人的劳动再生产出来。在生产过程中，工人创造的新价值不仅包括相当于劳动力价值的价值，还包括一定量的剩余价值($m$)。
\end{itemize}
基于此，提出资本主义商品价值构成公式：$W =c+v+m$. 将资本区分为不变资本和可变资本，进一步\textbf{解释了雇佣劳动者的剩余劳动是剩余价值的唯一源泉}。这种划分也为确定资本家对雇佣劳动者的剥削程度提供了科学依据。因此，要确定资本家对工人的剥削程度，就应该用剩余价值和雇佣劳动者的可变资本相比，即$m' = \frac{m}{v}$。
\subsection{剩余价值生产的两种基本方法}
\begin{itemize}
  \item 绝对剩余价值生产方法：在必要劳动时间不变的条件下，通过\textbf{延长工作时长和提高劳动强度}生产剩余价值。
  \item 相对剩余价值生产方法：通过\textbf{全社会劳动生产率的提高}实现。社会劳动生产率的提高，降低了劳动力的价值，从而缩短了必要劳动时间，相对延长了剩余劳动时间。（全社会劳动生产率的提高是资本家追逐超额剩余价值的结果；超额剩余价值是指企业由于提高劳动生产率而使商品的个别价值低于社会价值的差额。）
\end{itemize}
\subsection{资本积累}
\begin{mdquote}
把剩余价值转化为资本，或者说，\textbf{剩余价值的资本化}，就是资本积累。
\end{mdquote}
\begin{itemize}
  \item 本质：\textbf{资本积累是资本主义扩大再生产的源泉}，资本积累的本质就是资本家不断地利用无偿占有的工人创造的剩余价值来扩大自己的资本规模，进一步扩大和加强对工人的剥削与统治。
  \item 源泉：\textbf{资本积累的源泉是剩余价值}，资本积累规模的大小取决于资本家对工人的剥削程度、劳动生产率的高低、所用资本和所费资本之间的差额以及资本家预付资本的大小。
  \item 影响：\textbf{资本积累不但是社会财富占有两极分化的重要原因，而且是资本主义社会失业现象产生的根源。}
  \item 历史趋势：资本积累的历史趋势是资本主义制度的必然灭亡和社会主义制度的必然胜利。随着资本积累的不断增长，\textbf{生产社会化和生产资料资本主义私人占有}之间的深刻矛盾日益加剧，资本主义最终会被新的、更能适应社会化大生产要求的社会形态所取代。
\end{itemize}
\section{资本的循环、周转与再生产}
\subsection{资本循环和资本周转}
\begin{mdquote}
资本循环是资本从一种形式出发，经过一系列形式的变化，又回到原来出发点的运动。
\end{mdquote}
\begin{itemize}
  \item 产业资本在循环过程中要经历三个不同的阶段、执行三种不同的职能：
\begin{itemize}
  \item 第一个阶段是\textbf{购买}阶段，即生产资料与劳动力的购买阶段$\to$执行货币资本的职能；
  \item 第二个阶段是\textbf{生产}阶段，即生产资料与劳动力按比例结合在一起从事资本主义生产的阶段$\to$执行生产资本的职能；
  \item 第三个阶段是\textbf{售卖}阶段，即商品资本向货币资本的转化阶段$\to$执行商品资本的职能。
\end{itemize}
  \item 产业资本的运动必须具备两个前提条件：
\begin{itemize}
  \item 产业资本的三种职能形式必须\textbf{在空间上并存}；
  \item 产业资本的三种职能形式必须\textbf{在时间上继起}。
  \par 因此，产业资本的连续循环就是三种职能形式循环的统一。
\end{itemize}
\end{itemize}
\begin{mdquote}
资本是在运动中增殖的，必须不断地、周而复始地循环，才能不断地带来剩余价值。这种周而复始、不断反复着的资本循环，就叫作资本的周转。
\end{mdquote}
\begin{itemize}
  \item 影响资本周转快满的两个关键因素：
\begin{itemize}
  \item 一是资本周转的时间；
  \item 二是生产资本的固定资本和流动资本的构成。要加快资本周转速度，获得更多的剩余价值，就要缩短资本周转时间，加快流动资本的周转速度。
\end{itemize}
\end{itemize}
\subsection{社会再生产}
\begin{mdquote}
\begin{itemize}
  \item 社会生产是连续不断地进行的，这种连续不断重复的生产就是再生产。
  \item 所谓社会总产品就是社会在一定时期（通常为一年）所生产的全部物质资料的总和。
\end{itemize}
\end{mdquote}
\begin{itemize}
  \item 核心问题：社会再生产的核心问题是\textbf{社会总产品的实现问题}，即社会总产品的\textbf{价值补偿}和\textbf{实物补偿}问题。
\begin{itemize}
  \item 在价值形态上，社会总产品又叫社会总价值，划分为包括在产品中的生产资料的转移价值($c$)、凝结在产品中的由工人必要劳动时间创造的价值($v$)，以及凝结在产品中的由工人在剩余劳动时间里创造的价值($m$)。
  \item 在物质形态上，社会总产品根据其最终用途可以划分为用于生产消费的生产资料和用于生活消费的消费资料。
\end{itemize}
  \item 使社会总产品在价值上得到补偿，在实物上得到替换，以使社会再生产顺利进行，是不以人的意志为转移的客观规律。社会总产品在实物上得到替换，在价值上实现补偿，客观上就要求两大部类内部各个产业部门之间和两大部类之间保持一定的比例关系。
\begin{itemize}
  \item 生产资料的生产和消费资料的生产既要保证简单再生产的实现，更要保证扩大再生产的实现。
  \item 生产资料的生产既要满足两大部类对所耗费掉的生产资料的补偿，也要保证两大部类扩大生产规模后对追加生产资料的需求。
  \item 消费资料的生产则既要满足两大部类劳动者的个人消费和社会消费，也要满足扩大生产规模后对追加消费资料的需求。
  \item 上述比例不只是总量上的比例，还是结构上的比例。只有两大部类的生产不仅在规模上而且在结构上保持一定的比例，社会总产品的价值补偿和实物替换才能正常实现，社会再生产才能顺利进行。
\end{itemize}
\end{itemize}
\section{工资的本质}
\begin{mdquote}
在资本主义制度下，工人的工资表现为劳动力的价值或价格，这是资本主义工资的本质。
\end{mdquote}
\begin{itemize}
  \item 资本家购买工人的劳动力是以货币工资形式支付的，工人为资本家劳动，资本家付给工人工资，工资表现为劳动的价格或工人全部劳动的报酬，这就模糊了工人必要劳动和剩余劳动的界限，掩盖了资本主义的剥削关系。
  \item 在当代资本主义国家，随着社会生产力的发展、科学技术的进步、劳动过程的复杂化和脑力劳动作用的加强，工人的实际工资呈现出不断提高的趋势，但是与其创造的剩余价值的增长幅度相比，实际工资提高的幅度还是较小的。只要资本和雇佣劳动的基本经济关系不变，资本主义工资的本质就不会发生根本变化。
\end{itemize}
\section{资本主义的基本矛盾与经济危机}
\subsection{资本主义的基本矛盾}
\textbf{生产社会化和生产资料资本主义私人占有之间的矛盾}，是资本主义的基本矛盾。
\begin{itemize}
  \item 在资本主义条件下，随着科学技术的进步和社会生产力的不断提高，资本主义生产不断社会化。但是在资本家私人占有生产资料和剥削雇佣劳动者的生产关系中，\textbf{社会化的生产力却变成资本的生产力}，变成资本高效能地榨取剩余劳动、占有剩余价值、实现价值增殖的能力。这就形成了资本主义所特有的生产社会化和生产资料资本主义私人占有之间的矛盾。这是生产力和生产关系之间的矛盾在资本主义社会的具体体现。
  \item 资本主义越发展，科学技术以至社会生产力越发展，生产社会化的程度越高，不断发展的社会生产力就越成为资本的生产力，资本、生产资料、劳动产品就越来越集中在少数资本家的手里，资本主义基本矛盾的尖锐化就越是不可避免。
\end{itemize}
\subsection{资本主义经济危机}
\begin{mdquote}
资本主义发展到一定阶段，就会发生以\textbf{生产过剩}为基本特征的经济危机。
\end{mdquote}
\begin{itemize}
  \item 本质特征：生产过剩是资本主义经济危机的本质特征，但这种过剩是\textbf{相对过剩}，即\textbf{相对于劳动人民有支付能力的需求}来说社会生产的商品显得相对过剩，而不是与劳动人民的实际需要相比的绝对过剩。
  \item 形式原因：经济危机的抽象的一般的可能性，首先是由\textbf{货币作为流通手段和支付手段}引起的。
\begin{itemize}
  \item 以货币为媒介的商品买卖在时间上分为两个相互独立的行为。如果一些商品生产者在出卖自己的商品后不接着购买他人生产的商品，就会有另一些商品生产者的商品卖不出去；
  \item 同时，在商品买卖有更多的部分采取赊购赊销方式的情况下，如果某些债务人在债务到期时不能支付，就会是整个信用关系遭到破坏；但是，这仅仅是危机形式上的可能性。
\end{itemize}
  \item 根本原因：\textbf{资本主义的基本矛盾}，这一基本矛盾具体表现在以下两个方面：
\begin{itemize}
  \item 其一，生产无限扩大的趋势与劳动人民有支付能力的需求相对缩小的矛盾；
  \item 其二，单个企业内部生产的有组织性和整个社会生产的无政府状态之间的矛盾。
\end{itemize}
  \item 周期性：
\begin{itemize}
  \item 资本主义经济危机具有周期性，这是由\textbf{资本主义基本矛盾运动的阶段性}决定的。只要存在资本主义制度，经济危机就是不可避免的。
  \item 资本主义经济危机周期性爆发的特点，使社会资本再生产也呈现出周期性的特点，从一次危机开始到另一次危机的爆发，就是再生产的一个周期，一般包括危机、萧条、复苏和高涨四个相互联系的阶段。其中，危机阶段是必经阶段，没有危机阶段，就不存在资本主义再生产的周期性。
\end{itemize}
\end{itemize}
\chapter{第五章\quad{}资本主义的发展及其趋势}
\section{从自由竞争资本主义到垄断资本主义}
\subsection{垄断的形成}
\begin{mdquote}
垄断是指少数资本主义大企业为了获得高额利润，通过相互协议或联合，对一个或几个部门商品的生产、销售和价格进行操纵和控制。
\end{mdquote}
\begin{itemize}
  \item 产生的原因：
\begin{itemize}
  \item 当生产集中发展到相当高的程度，极少数企业就会联合起来，操纵和控制本部门的生产和销售，实行垄断，以获得高额利润；
  \item 企业规模巨大，形成对竞争的限制，也会产生垄断；
  \item 激烈的竞争给竞争各方带来的损失越来越严重，为了避免两败俱伤，企业之间会达成妥协，联合起来，实行垄断。
\end{itemize}
  \item 垄断的方式：
\begin{itemize}
  \item 垄断是通过一定的垄断组织形式实现的。垄断组织是指在一个或几个经济部门中，占据垄断地位的大企业联合。
  \item 尽管垄断组织的形式多种多样且不断变化发展，但本质上是一样的，即通过联合实现独占和瓜分商品生产和销售市场，操纵垄断价格，以攫取高额垄断利润。
\end{itemize}
\end{itemize}
\subsection{垄断下竞争的特点}
垄断作为自由竞争的对立面，并不能消除竞争，反而使竞争变得更加复杂和激烈，这是因为：
\begin{itemize}
  \item 垄断没有消除产生竞争的经济条件；
  \item 垄断必须通过竞争来维持；
  \item 社会生产是复杂多样的，任何垄断组织都不可能把包罗万象的社会生产全部包下来；
  \par 总之，在垄断条件下，在垄断组织内部、垄断组织之间以及垄断资本家集团之间，垄断组织同非垄断组织之间以及中小企业之间，存在广泛而激烈的竞争。
\end{itemize}
垄断条件下的竞争相比于自由竞争的新特点：\textbf{规模大、时间长、手段残酷、程度激烈，而且具有更大的破坏性}
\begin{center}
\begin{tabularx}{\linewidth}{lXX}
\toprule
 & \textbf{自由竞争} & \textbf{垄断} \\
\midrule
竞争目的 & 获取更多的利润或超额利润，不断扩大资本的积累 & \textbf{获取高额垄断利润}，并不断\textbf{巩固和扩大自己的垄断地位和统治权力} \\
竞争手段 & 主要运用经济手段，如通过改进技术、提高劳动生产率、降低产品成本 & 除了各种形式的经济手段之外，还采取\textbf{非经济的手段}，使竞争变得更加复杂激烈 \\
竞争范围 & 主要在经济领域、国内市场上进行 & \textbf{国际市场上的竞争越来越激烈}，不仅经济领域的竞争多种多样，而且还扩大到经济领域之外 \\
\bottomrule
\end{tabularx}
\end{center}
\subsection{金融资本与金融寡头}
\begin{itemize}
  \item 金融资本：由\textbf{工业垄断资本和银行垄断资本}融合在一起而形成的一种垄断资本。
  \item 金融寡头：在金融资本形成的基础上产生，指操纵国民经济命脉，并在实际上\textbf{控制国家政权的少数垄断资本家}或垄断资本家集团。他们支配了大量的社会财富，控制了整个国家的经济命脉和上层建筑，是垄断资本主义国家事实上的统治者。
\end{itemize}
\subsection{垄断利润和垄断价格}
\begin{mdquote}
垄断资本的实质在于获取垄断利润，垄断利润是垄断资本家凭借其在社会生产和流通中的垄断地位而获得的\textbf{超过平均利润}的高额利润。
\end{mdquote}
\begin{itemize}
  \item 垄断利润的来源：垄断资本所获得的高额利润，归根到底来\textbf{自无产阶级和其他劳动人民创造的剩余价值}。具体来说，垄断利润的来源大体有以下几个方面：
\begin{itemize}
  \item 来自对本国无产阶级和其他劳动人民剥削的加强；
  \item 由于垄断资本可以通过垄断高价和垄断低价来控制市场，使得它能获得一些其他企业特别是非垄断企业的利润；
  \item 通过加强对其他国家劳动人民的剥削和掠夺，从国外获取利润；
  \item 通过资本主义国家政权进行有利于垄断资本的再分配，从而将劳动人民创造的国民收入的一部分变成垄断资本的收入。
\end{itemize}
\end{itemize}
\begin{mdquote}
垄断利润主要是通过垄断组织指定的垄断价格来实现的。垄断价格是垄断组织在销售或购买商品时，凭借其垄断地位规定的、旨在保证获取最大限度利润的市场价格。
\end{mdquote}
\begin{itemize}
  \item 垄断价格的形式：垄断价格包括\textbf{垄断高价和垄断低价}两种形式。
\begin{itemize}
  \item 垄断高价是指垄断组织出售商品时规定的高于生产价格的价格；
  \item 垄断低价是指垄断组织在购买非垄断企业所生产的原材料等生产资料时规定的低于生产价格的价格。
\end{itemize}
  \item 垄断价格与价值规律：
\begin{itemize}
  \item 垄断组织操纵价格所带来的结果是抑制了市场上价格的自由波动，\textbf{垄断价格一定时期内背离生产价格和价值}。
  \item 但是，从全社会看，\textbf{整个社会商品的价值仍然是由生产它们的社会必要劳动时间决定的}，垄断价格既不能增加也不能减少整个社会所生产的价值总量，它只是对商品价值和剩余价值作了\textbf{有利于垄断资本的再分配}。
  \item 从全社会看，\textbf{商品的价格总额仍然等于商品的价值总额}。所以，垄断价格的产生没有否定价值规律，它是价值规律在垄断资本主义阶段作用的具体体现。
\end{itemize}
\end{itemize}
\section{垄断资本主义的发展}
\subsection{国家垄断资本主义的形成与作用}
\begin{mdquote}
国家垄断资本主义是\textbf{国家政权和私人垄断资本}融合在一起的垄断资本主义。国家垄断资本主义的产生，是垄断资本主义生产关系在自身范围内的部分质变，标志着资本主义发展进入了新的阶段。
\end{mdquote}
\begin{itemize}
  \item 成因：国家垄断资本主义的形成和发展不是偶然的，它是\textbf{科技进步和生产社会化程度进一步提高}的产物，是资本主义基本矛盾进一步尖锐化的必然结果：
\begin{itemize}
  \item \textbf{社会生产力的发展，要求资本主义生产资料在更大范围内被支配}，从而促进了国家垄断资本主义的产生；
  \item \textbf{经济波动和经济危机的深化}，要求国家垄断资本主义的产生；
  \item \textbf{缓和社会矛盾、协调利益关系}，要求国家垄断资本主义的产生。
\end{itemize}
  \item 五种主要形式：
\begin{itemize}
  \item \textbf{国家所有并直接经营企业}；
  \item \textbf{国家与私人共有、合营企业}；
  \item 国家通过多种形式\textbf{参与私人垄断资本的再生产过程}：
  \item \textbf{宏观调节}：国家运用\textbf{财政政策、货币政策}等经济手段，对社会总供给和总需求进行调节；
  \item \textbf{微观规制}：国家运用\textbf{法律手段}规范市场秩序，限制垄断，保护竞争，维护社会公众的合法权益：
\end{itemize}
  \item 主要作用：国家垄断资本主义是垄断资本主义的新发展，它对资本主义经济的发展产生了积极的作用。
\begin{itemize}
  \item 首先，国家垄断资本主义的出现在一定程度上\textbf{有利于社会生产力的发展，部分克服了社会化大生产和私人垄断资本之间的矛盾}；
  \item 其次，资产阶级国家凌驾于私人垄断资本之上，代表整个垄断资产阶级的利益，调节经济过程和经济活动，在一定范围内\textbf{突破了私人垄断资本的狭隘界限}；
  \item 再次，通过国家的\textbf{收入再分配手段，使劳动人民生活水平有所改善和提高}；
  \item 最后，在国家垄断资本主义的\textbf{参与和干预}下，各主要资本主义国家\textbf{现代化水平迅速提高}，国民经济现代化进程加快。
  \item 但是，国家垄断资本主义的出现并没有根本改变垄断资产阶级的性质，也没有从根本上消除资本主义的基本矛盾。
\end{itemize}
\end{itemize}
\subsection{垄断资本主义的实质}
\begin{itemize}
  \item 垄断资本主义的5个基本特征：
\begin{itemize}
  \item 垄断组织在经济生活中起决定作用；
  \item 在金融资本的基础上形成金融寡头的统治；
  \item 资本输出有了特别重要的意义；
  \item 瓜分世界的资本家国际垄断同盟已经形成；
  \item 最大资本主义大国已把世界上的领土瓜分完毕。
\end{itemize}
  \item 这些特征集中体现了帝国主义的实质，即垄断资本凭借垄断地位，获取高额垄断利润。为了获取高额垄断利润，垄断资本对内通过“参与制”和同政府的“个人联合＂谋求从经济到政治对整个国家的统治；对外运用经济、政治甚至战争的手段进行扩张，谋求对整个世界经济和政治的控制。
\end{itemize}
\section{经济全球化}
\begin{mdquote}
经济全球化是指在生产不断发展、科技加速进步、社会分工和国际分工不断深化、生产的社会化和国际化程度不断提高的情况下，世界各国、各地区的经济活动越来越超出某一国家和地区的范围而相互联系、相互依赖的过程。
\end{mdquote}
\begin{itemize}
  \item 表现（进程）：\textbf{生产全球化；贸易全球化；金融全球化}
  \item 动因：
\begin{itemize}
  \item 从本质上来讲，经济全球化是\textbf{生产力发展和社会化大生产的必然要求}：
  \item \textbf{科学技术的进步和生产力的发展}为经济全球化提供了坚实的物质基础和根本的推动力；
  \item \textbf{跨国公司的发展}为经济全球化提供了适宜的企业组织形式；
  \item \textbf{各国经济体制的变革和国际经济组织的发展}是经济全球化的体制与组织保障。
\end{itemize}
  \item 影响：
\begin{itemize}
  \item 经济全球化的过程是\textbf{生产社会化程度不断提高}的过程。在全球化进程中，社会分工得以在更大的范围内进行，\textbf{资金、技术等生产要素可以在国际社会流动和优化配置}，由此带来巨大的分工利益，推动世界生产力的发展。
  \item 一些发展中国家在参与经济全球化的进程中也得到了快速发展：
\begin{itemize}
  \item 经济全球化为发展中国家提供\textbf{先进技术和管理经验}；
  \item 经济全球化为发展中国家提供更多的\textbf{就业机会}；
  \item 经济全球化推动发展中国家\textbf{国际贸易发展}；
  \item 经济全球化促进发展中国家\textbf{跨国公司的发展}。
\end{itemize}
  \item 经济全球化也是一把“双刃剑”：
\begin{itemize}
  \item \textbf{发达国家和发展中国家在经济全球化过程中的地位和收益不平等、不平衡}；
  \item 加剧了发展中国家\textbf{资源短缺和环境污染}；
  \item 一定程度上增加\textbf{经济风险}。
\end{itemize}
\end{itemize}
\end{itemize}
\section{当代资本主义的新变化}
\begin{itemize}
  \item 表现：
\begin{itemize}
  \item 生产资料所有制的变化；
  \item 垄断资本形式的变化；
  \item 劳资关系和分配关系的变化；
  \item 社会阶层和阶级结构的变化；
  \item 经济调节机制和经济危机形态的变化；
  \item 政治制度的变化。
\end{itemize}
  \item 原因：
\begin{itemize}
  \item \textbf{科学技术革命和生产力的发展}，是当代资本主义发生新变化的根本推动力量；
  \item \textbf{工人阶级争取自身权利和利益的斗争}，是推动当代资本主义发生变化的重要力量；
  \item \textbf{社会主义制度初步显示的优越性}对当代资本主义产生了重要的影响；
  \item 主张改良主义的政党对资本主义制度的改革，也对当代资本主义新变化起到了重要作用。
\end{itemize}
  \item 实质：
\begin{itemize}
  \item 二战后资本主义发生的变化从根本上来说是人类社会发展的一般规律和资本主义经济规律作用的结果。二战后资本主义发生的变化是在资本主义制度基本框架内的变化，并不意味着资本主义生产关系的根本性质发生了变化；
  \item 当代资本主义虽然发生了一些新变化，但这种变化没有改变资本主义制度的本质，并没有克服资本主义的基本矛盾，也没有根本改变马克思主义关于资本主义的基本论断的科学性，根源于资本主义基本矛盾的经济危机依然是资本主义不可克服的痼疾。
  \item 正确认识资本主义的这些变化，有助于我们在深刻洞察资本主义本质的同时，实事求是地分析和借鉴资本主义发展过程中出现的符合社会化大生产要求的积极因素，为我所用，以进一步完善和发展社会主义制度。
\end{itemize}
\end{itemize}
\section{资本主义的历史地位和发展趋势}
\begin{itemize}
  \item 资本主义的历史地位：
\begin{itemize}
  \item 与封建社会相比，资本主义显示了巨大的\textbf{历史进步性}：
\begin{itemize}
  \item \textbf{资本主义将科学技术转变为强大的生产力}；
  \item 资本主义\textbf{追求剩余价值的内在动力和竞争的外在压力推动了社会生产力的迅速发展}；
  \item 资本主义的\textbf{意识形态和政治制度}作为上层建筑在战胜封建社会自给自足的小生产的生产方式，保护、促进和完善资本主义生产方式方面起着重要作用，从而推动了社会生产力的迅速发展，促进了社会进步。
\end{itemize}
  \item 资本主义的局限性：
\begin{itemize}
  \item \textbf{资本主义基本矛盾阻碍社会生产力的发展}；
  \item 资本主义制度下\textbf{财富占有两极分化，引起经济危机}；
  \item 资本家阶级支配和控制资本主义经济和政治的发展和运行，\textbf{不断激化社会矛盾和冲突}。
\end{itemize}
\end{itemize}
  \item 资本主义被社会主义所替代的历史必然性：\textbf{资本主义的内在矛盾决定了资本主义必然被社会主义所替代}。
\begin{itemize}
  \item \textbf{资本主义基本矛盾“包含着现代的一切冲突的萌芽”}；
\begin{itemize}
  \item 资本主义基本矛盾表现在\textbf{阶级关系上，是无产阶级和资产阶级的对立}；
  \item 资本主义基本矛盾表现在\textbf{生产上，是个别企业生产的有组织性和整个社会生产的无政府状态之间的对立}。
  \item 资本主义\textbf{经济危机的爆发}正是这个基本矛盾发展的结果。只有用社会主义生产方式取而代之，才能根本解决资本主义生产方式的基本矛盾。
\end{itemize}
  \item \textbf{资本积累推动资本主义基本矛盾不断激化并最终否定资本主义自身}；
\begin{itemize}
  \item 从资本积累过程来看，资本主义基本矛盾在基本积累过程中不断发展，资本的不断积累为否定资本主义制度自身准备了物质条件；
  \item 当资本主义基本矛盾及其派生的各种矛盾在资本积累中不断发展，资本主义生产关系不再适应生产力状况时，公有制取代私有制、社会主义取代资本主义就将成为不可避免的结果。这是资本积累过程所具有的客观历史趋势。
\end{itemize}
  \item \textbf{国家垄断资本主义是资本社会化的更高形式，将成为社会主义的前奏}；
\begin{itemize}
  \item 资本国有化将为社会主义革命提供直接的物质前提，是无产阶级社会主义革命的入口处。
  \item 到了国家垄断资本主义阶段，生产社会化、资本社会化和管理社会化都达到了资本主义生产方式的更高程度，从而为全社会共同占有生产资料和共同组织社会化生产准备了瞅给你发呢的物质条件和社会条件。
\end{itemize}
  \item \textbf{资本主义社会存在和资产阶级和无产阶级两大阶级之间的矛盾和斗争}。
\begin{itemize}
  \item 随着资本主义经济的发展，资产阶级由生产力的解放者日益变成阻碍者；
  \item 资本主义在造就了社会化大生产的同时，也产生了推动和运用这一先进生产力的无产阶级，在经济上所处的被剥削地位使无产阶级具有彻底的革命性和斗争精神。\textbf{社会化大生产使无产阶级成为最有组织性的革命力量}。
\end{itemize}
  \item 当今世界，虽然资本主义制度通过自我调节还能为生产力的发展提供一定的空间，但社会主义取代资本主义是历史的必然趋势。
\begin{itemize}
  \item 马克思、恩格斯关于资本主义社会基本矛盾的分析没有过时，关于资本主义必然消亡、社会主义必然胜利的历史唯物主义观点也没有过时，这是社会历史发展不可逆转的总趋势；
  \item 但道路是曲折的。资本主义最终消亡、社会主义最终胜利，必然是一个很长的历史过程。我们要深刻认识资本主义社会的自我调节能力，充分估计到西方发达国家在经济科技军事方面长期占据优势的客观现实，认真做好两种社会制度长期合作和斗争的各方面准备。
\end{itemize}
\end{itemize}
\end{itemize}
\chapter{第六章\quad{}社会主义的发展及其规律}
\section{社会主义从空想到科学的发展}
\begin{itemize}
  \item 随着社会化大生产的发展和资本主义生产方式的普遍确立，以及资本主义社会中生产社会化与生产资料资本主义私人占有之间的矛盾的激化，无产阶级与资产阶级的斗争更加激烈。无产阶级队伍不断壮大，并在与资产阶级的斗争中从自发走向自觉，表现出改造社会、创造历史的巨大力量。这些新的变化，为社会主义从空想发展为科学提供了社会需要和客观条件。
  \item 马克思恩格斯适应社会的需要，在新的历史条件下创立了\zhemph{唯物史观}和\zhemph{剩余价值学说}，为实现社会主义从空想到科学的飞跃奠定了坚实的理论基础。
  \item 马克思恩格斯在揭示人类社会发展一般规律和资本主义发展特殊规律的基础上，科学论证了科学社会主义代替资本主义的历史必然性，创立了科学社会主义学说，实现了社会主义从空想到科学的伟大飞跃。
\end{itemize}
\section{科学社会主义基本原则}
\begin{enumerate}
  \item 资本主义必然灭亡，社会主义必然胜利；
  \item 无产阶级是最先进最革命的阶级，肩负着推翻资本主义旧世界、建立社会主义和共产主义新世界的历史使命；
  \item 无产阶级革命是无产阶级进行斗争的最高形式，以建立无产阶级专政的国家为目的；
  \item 社会主义社会要在生产资料公有制基础上组织生产，以满足全体社会成员的需要为生产的根本目的；
  \item 社会主义社会要对社会主义生产进行有计划的指导和调节，实行按劳分配原则；
  \item 社会主义社会要合乎自然规律地改造和利用自然，努力实现人与自然的和谐共生；
  \item 社会主义社会必须坚持科学的理论指导，大力发展社会主义先进文化；
  \item 无产阶级政党是无产阶级的先锋队，社会主义事业必须始终坚持无产阶级政党的领导；
  \item 社会主义社会要大力解放和发展生产力，逐步消灭剥削和消除两极分化，实现共同富裕和社会全面进步，并最终向共产主义社会过渡；
  \item 共产主义是人类最美好的社会制度，实现共产主义是共产党人的最高理想。
\end{enumerate}
正确把握科学社会主义基本原则：
\begin{itemize}
  \item 必须始终坚持科学社会主义基本原则，反对任何背离科学社会主义一般原则的错误倾向；
  \item 要善于把科学社会主义基本原则与本国实际相结合，创造性地回答和解决社会主义革命、建设、改革中的重大问题；
  \item 紧跟时代和实践的发展，在不断总结新鲜经验中进一步丰富和发展科学社会主义基本原则。
\end{itemize}
\section{20世纪社会主义的探索及其经验}
\begin{itemize}
  \item 经济文化相对落后国家建设社会主义的长期性：
\begin{itemize}
  \item 生产力发展状况的制约；
  \item 经济基础和上层建筑发展状况的制约；
  \item 国际环境的严峻挑战；
  \item 马克思主义执政党对社会主义发展道路的探索和对社会主义建设规律的认识，需要一个长期的过程。
\end{itemize}
  \item 社会主义发展道路的多样性：
\begin{itemize}
  \item 各个国家的生产力的发展状况和社会发展阶段决定了社会主义发展道路具有不同的特点；
  \item 历史文化传统的差异性是造成不同国家社会主义发展道路多样性的重要条件；
  \item 时代和实践的不断发展，是造成社会主义发展道路多样性的现实原因。
\end{itemize}
  \item 探索符合本国国情的发展道路：
\begin{itemize}
  \item 探索社会主义发展道路，必须坚持对待马克思主义的科学态度；
  \item 探索社会主义发展道路，必须从当时当地的历史条件出发，坚持“走自己的路”；
  \item 探索社会主义发展道路，必须充分吸收人类一切文明成果。
\end{itemize}
  \item 社会主义在实践中开拓前进：
\begin{itemize}
  \item 在实践中开拓前进是社会主义事业发展的必然要求；
  \item 社会主义在实践中开拓前进必须遗循客观规律；
  \item 以自信担当、开拓奋进的姿态走向社会主义光明未来。
\end{itemize}
\end{itemize}
\chapter{第七章\quad{}共产主义崇高理想及其最终实现}
\section{马克思和恩格斯对共产主义的理解}
\begin{itemize}
  \item 物质财富极大丰富，消费资料按需分配；
  \item 社会关系高度和谐，人们精神境界极大提高；
  \item 实现每个人自由而全面的发展，人类从必然王国向自由王国飞跃。
\end{itemize}
\section{共产主义与人的自由全面发展}
实现人的自由而全面的发展，是马克思主义追求的根本价值目标，也是共产主义社会的根本特征。在共产主义社会，人的发展是自由而全面的发展，是建立在个体高度自由自觉的基础上的全面发展：
\begin{itemize}
  \item 首先，旧式分工的消除为人的自由全面发展创造了条件；
  \item 其次，自由时间大大延长，为人的自由全面发展提供了广阔的前景；
  \item 再次，在共产主义社会，劳动不再是单纯的谋生手段，而成为“生活的第一需要”；
  \item 最后，共产主义社会实现了人类解放，人类将最终从支配他们生活和命运的异己力量中解放出来，实现从必然王国向自由王国的飞跃，开始自觉地创造自己的历史。
\end{itemize}
\section{“两个必然”与“两个绝不会”的关系}
\begin{itemize}
  \item 两个必然：马克思恩格斯在共产党宣言中提出，资产阶级的灭亡和无产阶级的胜利是同样不可避免的，这就是我们常说的资本主义必然灭亡和社会主义必然胜利的两个必然。
  \item 两个绝不会：马克思在《\textless{}政治经济学批判\textgreater{}序言》中提出了两个绝不会，即无论哪一种社会形态，在他所容纳的全部生产力发挥出来之前是绝不会灭亡的；而新的更高的生产关系，在它的物质存在条件在旧社会的胎胞里成熟以前，是绝不会出现的。
  \item 关系：两个必然和两个绝不会是对资本主义的灭亡和共产主义的胜利的必然性以及这种必然性实现的时间和条件的全面论述，前者讲的是资本主义灭亡和共产主义胜利的客观必要性，是根本的方面；而后者讲的是这种必然性实现的时间和条件，它告诫我们，两个必然的实现需要客观条件，而在这个条件具备之前绝不会成为现实。
\end{itemize}
\section{共产主义远大理想与中国特色社会主义共同理想的关系}
\begin{itemize}
  \item 从时间上看，远大理想与共同理想的关系是最终理想与阶段性理想的关系；
  \item 从层次上来看，远大理想与共同理想的关系是最高纲领与最低纲领的关系；
  \item 从范围上来看，远大理想与共同理想的关系也是全人类理想与全体中国人民理想的关系。
  \item 辩证统一关系：必须以辩证思维把握和处理远大理想和共同理想的关系。任何时候都要坚持远大理想和共同理想的统一，不能把它们割裂开来、对立起来。没有远大理想的指引，就不会有共同理想的确立和坚持；没有共同理想的实现，远大理想就没有现实的基础。
\end{itemize}
\end{document}
